\exercisesection{Composition Product}

\begin{enumerate}
\def\labelenumi{\arabic{enumi})}
\tightlist
\item
  Let \(k\) be a commutative ring, A an associative unital \(k\) -algebra, M an (A, A)-bimodule. An extension of A by M is the data of an associative unital \(k\) -algebra B and of an exact sequence of \(k\) -modules
\end{enumerate}

\[
(\mathcal {E}): 0 \rightarrow \mathrm{M} \stackrel {{i}} {{\rightarrow}} \mathrm{B} \stackrel {{p}} {{\rightarrow}} \mathrm{A} \rightarrow 0,
\]

where p is an algebra homomorphism and where

\[
i (p (b) m) = b i (m) \quad \text { and } \quad i (m p (b)) = i (m) b \quad \text { for } \quad b \in \mathbf {B}, \quad m \in \mathbf {M}.
\]

Two extensions \(0 \to M \xrightarrow{i} B \xrightarrow{p} A \to 0\) and \(0 \to M \xrightarrow{i'} B' \xrightarrow{p'} A \to 0\) are said to be equivalent if there exists an algebra homomorphism \(u: B \to B'\) such that \(p' \circ u = p\) and \(u \circ i = i'\) ; the homomorphism \(u\) is then bijective.

\begin{enumerate}
\def\labelenumi{\alph{enumi})}
\item
  We assume that the exact sequence of \(k\) -modules \(0 \to \mathbf{M} \xrightarrow{i} \mathbf{B} \xrightarrow{p} \mathbf{A} \to 0\) is split; let \(s: \mathbf{A} \to \mathbf{B}\) be a \(k\) -linear section of \(p\) . Show that there exists an element \(f \in \mathbf{C}^2(\mathbf{A}, \mathbf{M})\) (X, p.~195, exercise 26, b)) such that \(s(a) s(a') - s(aa') = i(f(a, a'))\) for \(a, a' \in \mathbf{A}\) . Show that \(f\) is a 2-cocycle and that its class in \(\mathbf{H}^2(\mathbf{A}, \mathbf{M})\) is independent of the choice of \(s\) ; we denote it \(c(\mathcal{E})\) .
\item
  Let \(\mathrm{Ex}_{k}(\mathbf{A},\mathbf{M})\) be the set of equivalence classes of extensions of A by M which are trivial as extensions of k-modules. Show that c defines a bijection from \(\mathrm{Ex}_{k}(\mathbf{A},\mathbf{M})\) onto \(\mathrm{H}^{2}(\mathbf{A},\mathbf{M})\) .
\item
  Let \(\mathcal{E}:0\to\mathbf{M}\stackrel{i}{\to}\mathbf{B}\stackrel{p}{\to}\mathbf{A}\to0\) and \(\mathcal{E}':0\to\mathbf{M}\stackrel{i'}{\to}\mathbf{B}\stackrel{p'}{\to}\mathbf{A}\to0\) be two extensions of A by M. We denote by C the subalgebra of \(\mathbf{B}\times\mathbf{B}'\) consisting of the elements \((b,b')\) such that \(p(b)=p'(b')\) , c the ideal of C consisting of the elements \((i(m),-i'(m))\) for \(m\in\mathbf{M}\) , \(\mathbf{B}''\) the algebra C/c, \(\pi:\mathbf{C}\to\mathbf{B}''\) the quotient homomorphism. Let \(i'':\mathbf{M}\to\mathbf{B}''\) and \(p'':\mathbf{B}''\to\mathbf{A}\) be the homomorphisms such that \(i''(m)=\pi((m,0))\) for \(m\in\mathbf{M}\) and \(p''\pi((b,b'))=p(b)\) for \(b\in\mathbf{B},b'\in\mathbf{B}'\) . Show that the sequence \(0\to\mathbf{M}\to\mathbf{B}''\to\mathbf{A}\to0\) is an extension of A by M, called the sum extension of ( \(\mathcal{E}\) ) and ( \(\mathcal{E}'\) ), and that this defines a commutative group law on the set of equivalence classes of extensions of A by M. In particular, the set \(\mathrm{Ex}_{k}(\mathbf{A},\mathbf{M})\) is thereby endowed with a group structure, and the map \(c\) is a group isomorphism.
\end{enumerate}

\begin{enumerate}
\def\labelenumi{\arabic{enumi})}
\setcounter{enumi}{1}
\item
  Suppose that the \(k\) algebra A admits an augmentation \(\pi:\mathbf{A}\to k\) (X, p.~196, exercise 27) and that
\item
  Suppose that the \(k\) -algebra \(\mathbf{A}\) admits an augmentation \(\pi : \mathbf{A} \to k\) (X, p.~196, exercise 27) and that, moreover, \(\mathbf{A}\) is a projective \(k\) -module. We denote by I the kernel of \(\pi\) . Let \(\mathbf{M}\) be a left \(\mathbf{A}\) -module.
\end{enumerate}

\begin{enumerate}
\def\labelenumi{\alph{enumi})}
\item
  Let \((\mathcal{E}):0\to M_{(\pi)}\stackrel{i}{\to}\mathbf{B}\stackrel{p}{\to}\mathbf{A}\to 0\) be an extension of \(\mathbf{A}\) by the \((\mathbf{A},\mathbf{A})\) -bimodule \(M_{(\pi)}\) (X, p.~196, exercise 27); we denote by \(\overline{\mathbf{B}}\) the ideal of \(\mathbf{B}\) of elements \(b\in \mathbf{B}\) such that \(p(b)\in I\) . We have \(i(m)b = 0\) for \(m\in \mathbf{M}\) and \(b\in \overline{\mathbf{B}}\) , so that \(\overline{\mathbf{B}}\) is equipped with a structure of left \(\mathbf{A}\) -module such that \(p(\beta)b = \beta b\) for \(\beta \in \mathbf{B}\) , \(b\in \overline{\mathbf{B}}\) . Show that there exists an exact sequence of \(\mathbf{A}\) -modules \(0\to \mathrm{M}\to \overline{\mathrm{B}}\to I\to 0\) ; we denote by \(\theta (\mathcal{E})\) the class of this exact sequence in \(\operatorname {Ext}_{\mathbf{A}}^{1}(I,\mathbf{M})\) .
\item
  Show that \(\theta\) defines a map from \(\mathrm{Ex}_{k}(\mathbf{A},\mathbf{M}_{(\pi)})\) to \(\mathrm{Ext}_{\mathbf{A}}^{1}(I,\mathbf{M})\) . Prove the commutativity of the diagram:
\end{enumerate}

\[
\begin{array}{c c c} \operatorname{Ex} _ {k} (\mathbf {A}, \mathbf {M} _ {(\pi)}) & \stackrel {{\theta}} {{\longrightarrow}} & \operatorname{Ext} _ {\mathbf {A}} ^ {1} (I, \mathbf {M}) \\ \downarrow^ {c} & & \downarrow^ {- \delta} \\ \mathbf {H} ^ {2} (\mathbf {A}, \mathbf {M} _ {(\pi)}) & \stackrel {{\varphi}} {{\longrightarrow}} & \operatorname{Ext} _ {\mathbf {A}} ^ {2} (k, \mathbf {M}) \end{array}
\]

where δ is the connecting homomorphism deduced from the exact sequence \(0 \rightarrow I \rightarrow A \xrightarrow{\pi} k \rightarrow 0\) , and where φ is the isomorphism obtained by computing with the standard resolution (X, p.~196, exercise 27, b)). Deduce that θ is an isomorphism of groups.

\begin{enumerate}
\def\labelenumi{\arabic{enumi})}
\setcounter{enumi}{2}
\tightlist
\item
  Let G be a group, A the algebra \(\mathbf{Z}^{(\mathrm{G})}, j: \mathrm{G} \to \mathrm{A}\) the canonical injection. One endows A with the augmentation \(\pi: A \to Z\) defined by \(\pi(e_g) = 1\) for all \(g \in G\) . Let M be a left A-module.
\end{enumerate}

\begin{enumerate}
\def\labelenumi{\alph{enumi})}
\item
  Let \(0 \to M_{(\pi)} \xrightarrow{i} B \xrightarrow{p} A \to 0\) be an extension of \(A\) by the \((A, A)\) -bimodule \(M_{(\pi)}\) (X, p.~196, exercise 27). We denote by \(\Gamma\) the subset of \(B\) consisting of the elements \(b\) of \(B\) such that \(p(b) \in j(G)\) . Show that \(\Gamma\) , endowed with the multiplication induced by that of \(B\) , is a group, an extension of \(G\) by the \(G\) -module \(M\) .
\item
  The preceding construction defines a map from \(\mathrm{Ex}_{\mathbf{Z}}(\mathbf{A},\mathbf{M}_{(\pi)})\) to \(\mathrm{Ex}(\mathbf{G},\mathbf{M})\) (VIII, § 13). Show that the diagram
\end{enumerate}

\[
\begin{array}{c} \operatorname{Ex} _ {\mathbf {Z}} (\mathbf {A}, \mathbf {M} _ {(\pi)}) \to \operatorname{Ex} (\mathbf {G}, \mathbf {M}) \\ \downarrow^ {c} \qquad \qquad \qquad \qquad \qquad \qquad \qquad \qquad \qquad \qquad \downarrow \\ \mathbf {H} ^ {2} (\mathbf {A}, \mathbf {M} _ {(\pi)}) \xrightarrow {\varphi^ {2}} \mathbf {H} ^ {2} (\mathbf {G}, \mathbf {M}) \end{array}
\]

is commutative.

* 4) Let g be a Lie k-algebra, M a g-module. An extension of g by M is an exact sequence of k-modules \(0 \rightarrow M \xrightarrow{i} h \xrightarrow{p} g \rightarrow 0\) where h is a Lie k-algebra, p a homomorphism of Lie algebras, and where \(i(p(H)m) = [H, i(m)]\) for \(H \in h, m \in M\) . Two extensions

\[
0 \rightarrow \mathrm{M} \stackrel {i} {\rightarrow} \mathfrak {h} \stackrel {p} {\rightarrow} \mathfrak {g} \rightarrow 0, 0 \rightarrow \mathrm{M} \stackrel {i ^ {\prime}} {\rightarrow} \mathfrak {h} ^ {\prime} \stackrel {p ^ {\prime}} {\rightarrow} \mathfrak {g} \rightarrow 0
\]

are said to be equivalent if there exists a homomorphism \(u: \mathfrak{h} \to \mathfrak{h}'\) such that \(p' \circ u = p\) , \(u \circ i = i'\) . We denote by Liex (g, M) the set of equivalence classes of extensions of g by M. We assume that g is a free k-module.

\begin{enumerate}
\def\labelenumi{\alph{enumi})}
\tightlist
\item
  Let \((\mathcal{E}):0\to \mathrm{M}\stackrel {i}{\to}\mathfrak{h}\stackrel {p}{\to}\mathfrak{g}\to 0\) be an extension of \(\mathfrak{g}\) by \(\mathbf{M}\) ; we choose a \(k\) -linear section \(s:\mathfrak{g}\to \mathfrak{h}\) of \(p\) . We define an element \(f\in C^2 (\mathfrak{g},\mathbf{M})\) (X, p.~196, exercise 28, b)) by setting
\end{enumerate}

\[
i (f (x, y)) = [ s (x), s (y) ] - s ([ x , y ]) \quad \text { for } \quad x  ,   y \in \mathfrak {g}  .
\]

Show that \(f\) is a 2-cocycle and that its class in \(\mathbf{H}^{2}(\mathfrak{g},\mathbf{M})\) does not depend on the choice of \(s\) ; we denote it \(\theta(\mathcal{E})\) . b) Show that \(\theta\) defines a bijection from Liex \((\mathfrak{g},\mathbf{M})\) onto \(\mathbf{H}^{2}(\mathfrak{g},\mathbf{M})\) . Describe directly the group structure on Liex \((\mathfrak{g},\mathbf{M})\) obtained by transport via \(\gamma\) .

\begin{enumerate}
\def\labelenumi{\alph{enumi})}
\setcounter{enumi}{2}
\tightlist
\item
  Let U be the enveloping algebra of g, \(M_{(\pi)}\) the (U, U)-bimodule associated with M by means of the augmentation \(\pi: U \to k\) (X, p.~196, exercise 27). Let \(0 \to M_{(\pi)} \to B \xrightarrow{p} U \to 0\) be an extension of U by \(M_{(\pi)}\) (exercise 1); we denote by b the set of elements b of B such that \(p(b) \in j(g)\) .
\end{enumerate}

Show that b is a Lie k-algebra, an extension of g by M. This construction defines a map \(\lambda: \mathrm{Ex}_{k}(\mathrm{U}, \mathrm{M}_{(\pi)}) \to \mathrm{Liex}(\mathfrak{g}, \mathrm{M})\) . Show that the diagram

\[
\begin{array}{c c c} \operatorname{Ex} _ {k} (\mathrm{U}, \mathrm{M} _ {(\pi)}) & \xrightarrow {\lambda} & \text { Liex(g,M) } \\ \downarrow^ {c} & & \downarrow^ {\gamma} \\ \mathrm{H} ^ {2} (\mathrm{U}, \mathrm{M} _ {(\pi)}) & \xrightarrow {\varphi^ {2}} & \mathrm{H} ^ {2} (\mathrm{g,M}) \end{array}
\]

(where \(\varphi^2\) is the isomorphism defined in X, p.~196, exercise 27) is commutative. Deduce that \(\lambda\) is an isomorphism. *

¶ 5) Let G, F be two groups.

\begin{enumerate}
\def\labelenumi{\alph{enumi})}
\item
  For every extension \(\mathcal{E}: F \to E \to G\) of \(G\) by \(F\) (I, p.~62), the action of \(E\) on \(F\) by inner automorphisms defines a homomorphism \(E \to \text{Aut}(F)\) , whence, by passing to the quotient, a group homomorphism \(\theta: G \to \text{Aut}(F)/\text{Int}(F)\) , called the homomorphism associated with the extension \(\mathcal{E}\) . Two isomorphic extensions of \(G\) by \(F\) have the same associated homomorphism.
\item
  Conversely, we consider a group homomorphism \(\theta: G \to \text{Aut} (F)/\text{Int} (F)\) . Let
\end{enumerate}

\[
p: \operatorname{Aut} (\mathbf {F}) \rightarrow \operatorname{Aut} (\mathbf {F}) / \operatorname{Int} (\mathbf {F})
\]

the quotient map. We choose a map \(\sigma\) from G to Aut (F) such that \(p \circ \sigma = \theta\) , and a map \(f\) from G \(\times\) G to F such that \(\sigma(x)\, \sigma(y)\, (\sigma(xy))^{-1} = \text{Int}(f(x,y))\) for \(x,y \in G\) ; set \(c(x,y,z) = \sigma(x)\, (f(y,z)).f(x,yz).(f(xy,z))^{-1}.(f(x,y))^{-1}\) . Show that \(c(x,y,z)\) belongs to the center Z of F.

\begin{enumerate}
\def\labelenumi{\alph{enumi})}
\setcounter{enumi}{2}
\item
  We consider Z as a G-module by setting \(g . z = \sigma(g) (z)\) for \(g \in G, z \in Z\) ; this definition is independent of the choice of \(\sigma\) . The map c defines an element of \(C^{3}(G, Z)\) ; show that c is a 3-cocycle, and that its class in \(H^{3}(G, Z)\) is independent of the choices of \(\sigma\) and of f. We denote it \(\omega(G, F, \theta)\) . d) Show that the vanishing of \(\omega(G, F, \theta)\) is a necessary and sufficient condition for there to exist an extension of G by F having as associated homomorphism \(\theta\) . If \(\omega(G, F, \theta) = 0\) , show that one can choose f so that the cocycle c is zero, then define a group law on \(G \times F\) .
\item
  Assume \(\omega(G, F, \theta) = 0\) . Show that the set of isomorphism classes of extensions of \(G\) by \(F\) of the associated homomorphism \(\theta\) is a principal homogeneous set under the group H \(^{2}\) (G, Z).
\item
  Let H be a group, M a \(\mathbf{Z}^{(H)}\) -module, \(x\) be an element of \(\mathrm{H}^3 (\mathrm{H},\mathrm{M})\) . Show that there exists a group E with center M and a homomorphism \(\varphi :\mathrm{H}\to \mathrm{Aut}(\mathrm{E}) / \mathrm{Int}(\mathrm{E})\) , inducing on M the given H-module structure, such that \(\omega (\mathrm{H},\mathrm{E},\varphi) = x\) (take \(\mathbf{E} = \mathbf{M}\times \mathbf{L}\) , where L is the free group on \(\mathrm{H}\times \mathrm{H}\) ).
\end{enumerate}

\begin{enumerate}
\def\labelenumi{\arabic{enumi})}
\setcounter{enumi}{5}
\tightlist
\item
  Let M be a right A-module and N a left A-module. For \(n \geqslant 0\) , we denote by \(\mathcal{T}_{n}(M, N)\) the set of triples \((P, p, q)\) , where P is a complex of finitely generated free A-modules, such that \(P_{r} = 0\) for \(r < 0\) and \(r > n\) , and where \(p: P \to N\) and \(q: P^{*} \to M(n)\) are morphisms of A-complexes (we denote by \(P^{*}\) the Homgr complex \(_{A}\) \((P, A)\) ). Two elements \((P, p, q)\) and \((P', p', q')\) of \(\mathcal{T}_{n}(M, N)\) are said to be linked if there exists a morphism \(u: P \to P'\) such that \(p = p' \circ u\) and \(q' = q \circ 'u\) ; we denote by \(T_{n}(M, N)\) the quotient of \(\mathcal{T}_{n}(M, N)\) by the finest equivalence relation for which two related elements are equivalent (cf.~II, p.~52, exercise 9). a) Show that \(T_{n}(M, N)\) is identified with Tor \(_{n}^{A}\) \((M, N)\) .
\end{enumerate}

\begin{enumerate}
\def\labelenumi{\alph{enumi})}
\setcounter{enumi}{1}
\tightlist
\item
  Let \((\mathcal{E}):0\to \mathrm{M}^{\prime}\stackrel {i}{\to}\mathrm{M}\stackrel {\pi}{\to}\mathrm{M}^{\prime \prime}\to 0\) be an exact sequence of right A-modules. Describe the connecting homomorphism \(\partial :\mathbf{T}_n(\mathbf{M}'',\mathbf{N})\to \mathbf{T}_{n - 1}(\mathbf{M}',\mathbf{N})\) deduced from \(\partial (\mathcal{E},\mathbf{N})\) using the previous identifications.
\end{enumerate}

\((\mathrm{If}(\mathbf{P},p,q)\in\mathcal{T}_{n}(\mathbf{M}^{\prime\prime},\mathbf{N}),\text{show that we have}\partial(\mathbf{P},p,q)=(\overline{\mathbf{P}},\overline{p},\overline{q}),\text{where}\overline{\mathbf{P}}\text{ is the complex obtained by replacing } \mathbf{P}_{n}\text{ by 0 in } \mathbf{P},\overline{p}\text{ is the morphism deduced from } p\text{ and }(\overline{q})^{n-1}: \mathbf{P}_{n-1}^{*}\to\mathbf{M}^{\prime}\text{ is equal to }u^{n-1},u\text{ being a morphism from } \mathbf{P}^{*}\text{ into the complex (Con (i)) (n) such that }\pi\circ u^{n}=q^{n}.)\)

¶ 7) Let A' be an (associative and unital) k-algebra and B = A ⊗k A'.

\begin{enumerate}
\def\labelenumi{\alph{enumi})}
\tightlist
\item
  Let M be a left A-module, Q a right A-module, M' a left A'-module, Q' a right A'-module. Let (P, p) (resp. (P', p'), resp. (R, r)) be a projective resolution of the A-module M (resp. of the A'-module M', resp. of the B-module M ⊗\_k M'). Show that there exists a morphism of B-complexes
\end{enumerate}

\[
u: \mathbf {P} \otimes_ {k} \mathbf {P} ^ {\prime} \rightarrow \mathbf {R},
\]

unique up to homotopy, such that \(r \circ u = p \otimes p'\) . Define, using \(u\) , a graded \(k\) -homomorphism of degree zero

\[
\top : \operatorname{Tor} ^ {\mathbf {A}} (\mathrm{Q}, \mathrm{M}) \otimes_ {k} \operatorname{Tor} ^ {\mathbf {A} ^ {\prime}} (\mathrm{Q} ^ {\prime}, \mathrm{M} ^ {\prime}) \rightarrow \operatorname{Tor} ^ {\mathbf {B}} (\mathrm{Q} \otimes_ {k} \mathrm{Q} ^ {\prime}, \mathrm{M} \otimes_ {k} \mathrm{M} ^ {\prime}).
\]

Show that \(\top\) does not depend on the choice of resolutions P, P', R and of the morphism \(u\) . If \(a \in \operatorname{Tor}^{\mathbf{A}}(\mathbf{Q}, \mathbf{M})\) and \(a' \in \operatorname{Tor}^{\mathbf{A}'}(\mathbf{Q}', \mathbf{M}')\) , we set \(\top(a \otimes a') = a \top a'\) .

\begin{enumerate}
\def\labelenumi{\alph{enumi})}
\setcounter{enumi}{1}
\tightlist
\item
  We assume moreover that the \(k\) -modules \(\mathbf{A}\) and \(\mathbf{A}'\) are projective and that one has \(\operatorname{Tor}_n^k (\mathbf{M},\mathbf{M}') = 0\) for \(n > 0\) . Show that \(\mathbf{P} \otimes_k \mathbf{P}'\) is then a projective resolution of the \(\mathbf{B}\) -module \(\mathbf{M} \otimes_k \mathbf{M}'\) ; deduce from this, for every \(\mathbf{A}\) -module \(\mathbf{N}\) and for every \(\mathbf{A}'\) -module \(\mathbf{N}'\) , a graded \(k\) -homomorphism of degree zero
\end{enumerate}

\[
\vee : \operatorname{Ext} _ {\mathbf {A}} (\mathbf {M}, \mathbf {N}) \otimes_ {k} \operatorname{Ext} _ {\mathbf {A} ^ {\prime}} (\mathbf {M} ^ {\prime}, \mathbf {N} ^ {\prime}) \rightarrow \operatorname{Ext} _ {\mathbf {B}} (\mathbf {M} \otimes_ {k} \mathbf {M} ^ {\prime}, \mathbf {N} \otimes_ {k} \mathbf {N} ^ {\prime}).
\]

Show that \(\vee\) is independent of the choice of resolutions \(\mathbf{P}\) and \(\mathbf{P}'\) . If \(b \in \operatorname{Ext}_{\mathbf{A}}(\mathbf{M}, \mathbf{N})\) and \(b' \in \operatorname{Ext}_{\mathbf{A}'}(\mathbf{M}', \mathbf{N}')\) , we write \(\vee (b \otimes b') = b \vee b'\) .

\begin{enumerate}
\def\labelenumi{\alph{enumi})}
\setcounter{enumi}{2}
\tightlist
\item
  Let A'' be a third k-algebra (associative and unital), M'' and N'' left A''-modules, Q'' a right A''-module. We identify the k-algebras A ⊗\_k (A' ⊗\_k A'') and (A ⊗\_k A') ⊗\_k A'', as well as the k-modules M ⊗\_k (M' ⊗\_k M'') and (M ⊗\_k M') ⊗\_k M'', N ⊗\_k (N' ⊗\_k N'') and (N ⊗\_k N') ⊗\_k N'', Q ⊗\_k (Q' ⊗\_k Q'') and (Q ⊗\_k Q') ⊗\_k Q''. Prove the formulas
\end{enumerate}

\[
\begin{array}{l l} a \top (a ^ {\prime} \top a ^ {\prime \prime}) = (a \top a ^ {\prime}) \top a ^ {\prime \prime} & \text {for} \quad a \in \operatorname{Tor} ^ {\mathbf {A}} (\mathrm{Q}, \mathrm{M}), \quad a ^ {\prime} \in \operatorname{Tor} ^ {\mathbf {A} ^ {\prime}} (\mathrm{Q} ^ {\prime}, \mathrm{M} ^ {\prime}), \\ & \qquad \qquad \qquad \qquad \qquad \qquad \qquad \qquad \qquad \qquad \qquad \qquad \qquad \qquad a ^ {\prime \prime} \in \operatorname{Tor} ^ {\mathbf {A} ^ {\prime \prime}} (\mathrm{Q} ^ {\prime \prime}, \mathrm{M} ^ {\prime \prime}) \end{array}
\]

\[
\begin{array}{r l} b \vee (b ^ {\prime} \vee b ^ {\prime \prime}) = (b \vee b ^ {\prime}) \vee b ^ {\prime \prime} & \text {for} \quad b \in \operatorname{Ext} _ {\mathbf {A}} (\mathbf {M}, \mathbf {N}), \quad b ^ {\prime} \in \operatorname{Ext} _ {\mathbf {A}} (\mathbf {M} ^ {\prime}, \mathbf {N} ^ {\prime}), \\ & \qquad \qquad \qquad \qquad \qquad \qquad b ^ {\prime \prime} \in \operatorname{Ext} _ {\mathbf {A} ^ {\prime \prime}} (\mathbf {M} ^ {\prime \prime}, \mathbf {N} ^ {\prime \prime}). \end{array}
\]

\begin{enumerate}
\def\labelenumi{\alph{enumi})}
\setcounter{enumi}{3}
\tightlist
\item
  Let \(\sigma : \mathrm{Tor}^{\mathbf{A} \otimes_{\mathbf{k}} \mathbf{A}'}(\mathbf{Q} \otimes_{\mathbf{k}} \mathbf{Q}', \mathbf{M} \otimes_{\mathbf{k}} \mathbf{M}') \to \mathrm{Tor}^{\mathbf{A}' \otimes_{\mathbf{k}} \mathbf{A}}(\mathbf{Q}' \otimes_{\mathbf{k}} \mathbf{Q}, \mathbf{M}' \otimes_{\mathbf{k}} \mathbf{M})\) and
\end{enumerate}

\[
\tau : \operatorname{Ext} _ {\mathbf {A} \otimes_ {k} \mathbf {A} ^ {\prime}} (\mathrm{M} \otimes_ {k} \mathrm{M} ^ {\prime}, \mathrm{N} \otimes_ {k} \mathrm{N} ^ {\prime}) \rightarrow \operatorname{Ext} _ {\mathbf {A} ^ {\prime} \otimes_ {k} \mathbf {A}} (\mathrm{M} ^ {\prime} \otimes_ {k} \mathrm{M}, \mathrm{N} ^ {\prime} \otimes_ {k} \mathrm{N})
\]

the isomorphisms deduced from the commutation isomorphisms

\[
\mathrm{A} \otimes_ {k} \mathrm{A} ^ {\prime} \rightarrow \mathrm{A} ^ {\prime} \otimes_ {k} \mathrm{A}, \quad \mathrm{M} \otimes_ {k} \mathrm{M} ^ {\prime} \rightarrow \mathrm{M} ^ {\prime} \otimes_ {k} \mathrm{M}, \quad \mathrm{N} \otimes_ {k} \mathrm{N} ^ {\prime} \rightarrow \mathrm{N} ^ {\prime} \otimes_ {k} \mathrm{N}, \quad \mathrm{Q} \otimes_ {k} \mathrm{Q} ^ {\prime} \rightarrow \mathrm{Q} ^ {\prime} \otimes_ {k} \mathrm{Q}.
\]

Prove the formulas.

\[
\begin{array}{l l} \sigma (a \top a ^ {\prime}) = (- 1) ^ {p q} a ^ {\prime} \top a & \text { for } \quad a \in \operatorname{Tor} _ {p} ^ {A} (Q, M), \quad a ^ {\prime} \in \operatorname{Tor} _ {q} ^ {A ^ {\prime}} (Q ^ {\prime}, M ^ {\prime}) \\ \tau (b \vee b ^ {\prime}) = (- 1) ^ {p q} b ^ {\prime} \vee b & \text { for } \quad b \in \operatorname{Ext} _ {A} ^ {p} (M, N), \quad b ^ {\prime} \in \operatorname{Ext} _ {A ^ {\prime}} ^ {q} (M ^ {\prime}, N ^ {\prime}). \end{array}
\]

\begin{enumerate}
\def\labelenumi{\alph{enumi})}
\setcounter{enumi}{4}
\tightlist
\item
  Let \((\mathcal{E}):0\to \mathbf{M}_2\to \mathbf{M}_1\to \mathbf{M}_0\to 0\) be an exact sequence of A-modules, such that the sequence
\end{enumerate}

\[
(\mathcal {E} \otimes \mathrm{M} ^ {\prime}): 0 \rightarrow \mathrm{M} _ {2} \otimes_ {k} \mathrm{M} ^ {\prime} \rightarrow \mathrm{M} _ {1} \otimes_ {k} \mathrm{M} ^ {\prime} \rightarrow \mathrm{M} _ {0} \otimes_ {k} \mathrm{M} ^ {\prime} \rightarrow 0
\]

be exact. Let δ (resp. Δ) denote the connecting homomorphism δ(Q, ℰ) (resp. δ(Q ⊗ₖ Q′, ℰ ⊗ M′)). Show that one has (δa) ⊤ a′ = Δ(a ⊗ a′) for a ∈ Tor\^{}A(Q, M₀), a′ ∈ Tor\^{}A(Q', M'). Prove the analogous formulas for the product ∨.

\begin{enumerate}
\def\labelenumi{\alph{enumi})}
\setcounter{enumi}{5}
\tightlist
\item
  If k is a field, show that \(\top\) is an isomorphism of graded k-modules; the same holds for ∨ if, moreover, A and A' are left noetherian and the A-module M and the A'-module M' are finitely generated. g) If one identifies the modules \(\operatorname{Tor}_{n}^{\mathbf{A}}(\mathbf{Q},\mathbf{M})\) with the sets \(\operatorname{T}_{n}(\mathbf{Q},\mathbf{M})\) defined in Exercise 6, prove the formula \((\mathbf{P},p,q)\top(\mathbf{P}',p',q')=(\mathbf{P}\otimes_{k}\mathbf{P}',p\otimes p',q\otimes q')\) .
\end{enumerate}

\begin{enumerate}
\def\labelenumi{\arabic{enumi})}
\setcounter{enumi}{7}
\tightlist
\item
  We keep the notation of Exercise 7; we assume that the \(k\) -modules A and A' are projective.
\end{enumerate}

\begin{enumerate}
\def\labelenumi{\alph{enumi})}
\item
  We assume that one has \(\mathrm{Tor}_{i}^{k}(\mathbf{M},\mathbf{M}^{\prime}) = 0\) for \(i > 0\) . Show that the map \(b\mapsto b\vee 1_{\mathbf{M}'}\) defines a graded \(k\) -homomorphism of degree zero from \(\operatorname{Ext}_{\mathbf{A}}(\mathbf{M},\mathbf{N})\) to \(\operatorname{Ext}_{\mathbf{B}}(\mathbf{M}\otimes_{k}\mathbf{M}',\mathbf{N}\otimes_{k}\mathbf{M}')\) , which induces in degree zero the canonical homomorphism \(b\mapsto b\otimes 1_{\mathbf{M}'}\) .
\item
  We assume that the k-module M' is flat. Let
\end{enumerate}

\[
0 \rightarrow \mathrm{N} \rightarrow \mathrm{R} _ {n} \rightarrow \dots \rightarrow \mathrm{R} _ {1} \rightarrow \mathrm{M} \rightarrow 0
\]

be an exact sequence of A-modules, of class \(b \in \operatorname{Ext}_{\mathbf{A}}^{n}(\mathbf{M}, \mathbf{N})\) . Show that the exact sequence of B-modules

\[
0 \rightarrow \mathrm{N} \otimes_ {k} \mathrm{M} ^ {\prime} \rightarrow \mathrm{R} _ {n} \otimes_ {k} \mathrm{M} ^ {\prime} \rightarrow \dots \rightarrow \mathrm{R} _ {1} \otimes_ {k} \mathrm{M} ^ {\prime} \rightarrow \mathrm{M} \otimes_ {k} \mathrm{M} ^ {\prime} \rightarrow 0
\]

has class \(b \vee 1_{\mathbf{M}'}\) in \(\operatorname{Ext}_{\mathbf{B}}^{n}(\mathbf{M} \otimes_{k} \mathbf{M}', \mathbf{N} \otimes_{k} \mathbf{N}')\) .

\begin{enumerate}
\def\labelenumi{\alph{enumi})}
\setcounter{enumi}{2}
\tightlist
\item
  We assume that the k-modules M, M', N, N' are flat. Show that we have:
\end{enumerate}

\[
b \vee b ^ {\prime} = (b \vee 1 _ {N ^ {\prime}}) \circ (1 _ {M} \vee b ^ {\prime}) = (- 1) ^ {p q} (1 _ {N} \vee b ^ {\prime}) \circ (b \vee 1 _ {M ^ {\prime}})
\]

for \(b \in \operatorname{Ext}_{\mathbf{A}}^{p}(\mathbf{M}, \mathbf{N})\) , \(b' \in \operatorname{Ext}_{\mathbf{A}'}^{q}(\mathbf{M}', \mathbf{N}')\) .

\begin{enumerate}
\def\labelenumi{\arabic{enumi})}
\setcounter{enumi}{8}
\tightlist
\item
  We assume the ring A is commutative. Let B, C be two associative and unital A-algebras; we denote by \(m_{\mathbf{B}}: \mathbf{B} \otimes_{\mathbf{A}} \mathbf{B} \to \mathbf{B}\) the multiplication of B, and \(m_{\mathbf{C}}\) the multiplication of C.
\end{enumerate}

\begin{enumerate}
\def\labelenumi{\alph{enumi})}
\tightlist
\item
  Show that the homomorphism
\end{enumerate}

\[
\operatorname{Tor} \left(m _ {\mathrm{B}}, m _ {\mathrm{C}}\right) \circ \top : \operatorname{Tor} ^ {\mathrm{A}} (\mathrm{B}, \mathrm{C}) \otimes_ {\mathrm{A}} \operatorname{Tor} ^ {\mathrm{A}} (\mathrm{B}, \mathrm{C}) \rightarrow \operatorname{Tor} ^ {\mathrm{A}} (\mathrm{B}, \mathrm{C})
\]

(cf.~exercise 7) defines on Tor \(^{A}\) (B, C) a structure of graded A-algebra, associative and unital. b) Let S be a differential graded A-algebra and \(p : S_{0} \to B\) a homomorphism of A-algebras such that (S, p) is a free resolution of B (cf. X, p. 183, exercise 18). Show that the isomorphism

\[
\psi (\mathbf {S}, \mathbf {C}): \operatorname{Tor} ^ {\mathbf {A}} (\mathbf {B}, \mathbf {C}) \rightarrow \mathbf {H} (\mathbf {S} \otimes_ {\mathbf {A}} \mathbf {C})
\]

is an isomorphism of A-algebras if one endows H(S ⊗A C) with the algebra structure deduced from the differential graded algebra structure of S ⊗A C (X, p.~183, exercise 18, b) and c)).

\begin{enumerate}
\def\labelenumi{\alph{enumi})}
\setcounter{enumi}{2}
\tightlist
\item
  If the A-algebras B and C are commutative, prove that the graded algebra Tor \(^{A}\) (B, C) is alternating (cf. X, p. 183, exercise 19).
\end{enumerate}

\begin{enumerate}
\def\labelenumi{\arabic{enumi})}
\setcounter{enumi}{9}
\tightlist
\item
  Let B be a k-bialgebra (III, p.~148), which we regard as an augmented k-algebra (X, p.~196, exercise 27) via the counit \(\beta: \mathbf{B} \to k\) ; we assume that the \(k\) -module B is projective.
\end{enumerate}

\begin{enumerate}
\def\labelenumi{\alph{enumi})}
\tightlist
\item
  Let M, N be two left B-modules. The homomorphism ∨ of Exercise 7 is identified with a graded homomorphism of degree zero
\end{enumerate}

\[
\mathrm{H} ^ {\bullet} (\mathbf {B}, \gamma ; \mathbf {M}) \otimes_ {k} \mathrm{H} ^ {\bullet} (\mathbf {B}, \gamma ; \mathbf {N}) \rightarrow \mathrm{H} ^ {\bullet} (\mathbf {B} \otimes_ {k} \mathbf {B}, \gamma \otimes \gamma ; \mathbf {M} \otimes_ {k} \mathbf {N});
\]

by composing with the homomorphism induced by the comultiplication, deduce from this a graded homomorphism of degree zero \(\cup: H'(B, \gamma; M) \otimes_k H'(B, \gamma; N) \to H'(B, \gamma; M \otimes_k N)\) .

\begin{enumerate}
\def\labelenumi{\alph{enumi})}
\setcounter{enumi}{1}
\item
  Let C be an associative unital k-algebra, which we equip with the B-module structure induced by γ. Show that the homomorphism ∪ defines on H'(B, γ ; C) a structure of associative unital graded k-algebra; if C is commutative and the bialgebra B is cocommutative, the algebra H'(B, γ ; C) is anticommutative.
\item
  We take C = k. Show that the product ∪ coincides with the composition product on \(\operatorname{Ext}_{\mathbf{B}}(k, k)\) . d) Let G be a group, M and N two \(\mathbf{Z}^{(\mathbf{G})}\) -modules. We endow M ⊗\_z N with the structure of \(\mathbf{Z}^{(\mathbf{G})}\) -module defined by \(g(x \otimes y) = gx \otimes gy\) for \(g \in G\) , \(x \in M\) , \(y \in N\) . Show that the map
\end{enumerate}

\[
\cup : \mathrm{H} ^ {\bullet} (\mathrm{G}, \mathrm{M}) \otimes_ {\mathbf {z}} \mathrm{H} ^ {\bullet} (\mathrm{G}, \mathrm{N}) \rightarrow \mathrm{H} ^ {\bullet} (\mathrm{G}, \mathrm{M} \otimes_ {\mathbf {z}} \mathrm{N})
\]

is the homomorphism induced on homology by the graded homomorphism of degree zero

\[
u: \mathrm{C} ^ {*} (\mathrm{G}, \mathrm{M}) \otimes_ {\mathbf {z}} \mathrm{C} ^ {*} (\mathrm{G}, \mathrm{N}) \rightarrow \mathrm{C} ^ {*} (\mathrm{G}, \mathrm{M} \otimes_ {\mathbf {z}} \mathrm{N})
\]

defined by \(u(f \otimes h)(g_1, \ldots, g_{p+q}) = f(g_1, \ldots, g_p) \otimes g_1 \ldots g_p h(g_{p+1}, \ldots, g_{p+q})\) for

\[
f \in \mathbf {C} ^ {p} (\mathbf {G}, \mathbf {M}), \quad h \in \mathbf {C} ^ {q} (\mathbf {G}, \mathbf {N}), \quad g _ {1} \dots g _ {p + q} \in \mathbf {G}.
\]

(Use Exercise 5, p.~185.)

¶ 11) Suppose that the bigebra B admits an inversion i (cf.~III, p.~198, exercise 4). Let M and N be B-modules. The k-module M ⊗k N (resp. Homk(M, N)) has a natural structure of B ⊗k B-module (resp. B ⊗k B°-module); one endows it with the B-module structure deduced from the homomorphism c (resp. (1B ⊗ i) ∘ c).

\begin{enumerate}
\def\labelenumi{\alph{enumi})}
\tightlist
\item
  Show that the homomorphism u (Exercise 10) defines a graded homomorphism of degree zero
\end{enumerate}

\[
\cup : \mathrm{H} ^ {*} (\mathbf {B}, \gamma ; \operatorname{Hom} _ {k} (\mathbf {M}, \mathbf {N})) \otimes_ {k} \mathrm{H} ^ {*} (\mathbf {B}, \gamma ; \mathbf {M}) \rightarrow \mathrm{H} ^ {*} (\mathbf {B}, \gamma ; \mathbf {N}).
\]

\begin{enumerate}
\def\labelenumi{\alph{enumi})}
\setcounter{enumi}{1}
\tightlist
\item
  Let \(0 \to M \to R_n \to \cdots \to R_1 \to k \to 0\) be an exact sequence of B-modules, with \(R_i\) projective ( \(1 \leqslant i \leqslant n\) ); let \(\theta\) be its class in \(\operatorname{Ext}_B^n(k, M) = H^n(B, \gamma; M)\) . Show that for \(p \geqslant 1\) , the map \(x \mapsto x \cup \theta\) is a \(k\) -isomorphism of \(H^p(B, \gamma; \operatorname{Hom}_k(M, N))\) onto \(H^{p+n}(B, \gamma; N)\) .
\end{enumerate}

(Show that this map factors through

\[
\mathrm{H} ^ {p} (\mathrm{B}, \gamma ; \operatorname{Hom} _ {k} (\mathrm{M}, \mathrm{N})) \rightarrow \operatorname{Ext} _ {\mathrm{B}} ^ {p} (\mathrm{M}, \mathrm{N}) \xrightarrow {\theta} \operatorname{Ext} _ {\mathrm{B}} ^ {p + n} (k, \mathrm{N}).
\]

\begin{enumerate}
\def\labelenumi{\alph{enumi})}
\setcounter{enumi}{2}
\tightlist
\item
  Let G be a group, N a \(\mathbf{Z}^{(G)}\) -module. Let F be a free group, \(p: F \to G\) a surjective homomorphism, R = Ker( \(p\) ), \(\overline{R} = R/(R, R)\) , \(\overline{F} = F/(R, R)\) . The extension \(\overline{R} \to \overline{F} \to G\) defines a class \(\varepsilon \in H^2(G, \overline{R})\) ; prove that the map \(x \mapsto x \cup \varepsilon\) from \(H^p(G, \operatorname{Hom}_{\mathbf{Z}}(\overline{R}, N))\) to \(H^{p+2}(G, N)\) is an isomorphism for \(p \geqslant 1\) (use exercise 3, d), p.~179).
\end{enumerate}

\exercisesection{Homological Dimension}

\begin{enumerate}
\def\labelenumi{\arabic{enumi})}
\tightlist
\item
  Let \(u: \mathbf{A} \to \mathbf{B}\) be a ring homomorphism, M a B-module. Prove the inequality
\end{enumerate}

\[
\mathrm{dp} _ {\mathbf {A}} (\mathbf {M}) \leqslant \mathrm{dp} _ {\mathbf {B}} (\mathbf {M}) + \mathrm{dp} _ {\mathbf {A}} (\mathbf {B} _ {s}).
\]

(Reason by induction on dpB(M) or use the spectral sequence of Exercise 6, p.~188.)

\begin{enumerate}
\def\labelenumi{\arabic{enumi})}
\setcounter{enumi}{1}
\item
  Assume the ring A is noetherian; let M be a finitely generated A-module of projective dimension n. Show that the k-module Ext \(_{A}^{n}\) (M, A) is nonzero.
\item
  Let \(n\) be an integer \(\geqslant 1\) . Show that \(\mathrm{dh}(\mathbf{M}_n(\mathbf{A})) = \mathrm{dh}(\mathbf{A})\) (cf.~VIII, § 4).
\item
  Let M be an A-module. One calls the injective dimension of M, and denotes it by \(\mathrm{di}_{\mathbf{A}}(\mathbf{M})\) , the infimum in \(\overline{\mathbf{Z}}\) of the lengths of injective resolutions of M.
\end{enumerate}

If n is an integer \(\geqslant\) 0, show that the following conditions are equivalent:

\(\alpha)\mathrm{di}_{\mathbf{A}}(\mathbf{M})\leqslant n;\)

β) \(\operatorname{Ext}_{\mathbf{A}}^{r}(\mathbf{N},\mathbf{M}) = 0\) for every A-module N and every integer \(r > n\) ;

\(\gamma)\ \mathrm{Ext}_{\mathbf{A}}^{n+1}(\mathbf{N},\mathbf{M}) = 0\) for every A-module \(\mathbf{N}\) ;

δ) \(\operatorname{Ext}_{\mathbf{A}}^{n+1}(\mathbf{N},\mathbf{M})=0\) for every cyclic A-module N;

ε) for every exact sequence 0 → M → I⁰ → \(\cdots\) → Iⁿ⁻¹ → N → 0, where the Iᵏ are injective, N is injective. 5) Let (Mᵢ)ᵢ∈F be a family of A-modules.

\begin{enumerate}
\def\labelenumi{\alph{enumi})}
\item
  Show that \(\mathrm{di}_{\mathbf{A}}(\prod \mathbf{M}_i) = \sup \mathrm{di}_{\mathbf{A}}(\mathbf{M}_i)\) .
\item
  If A is noetherian, show that \(\mathrm{di}_{\mathbf{A}}(\bigoplus \mathbf{M}_i) = \sup \mathrm{di}_{\mathbf{A}}(\mathbf{M}_i)\) .
\item
  If \(\mathrm{di}_{\mathbf{A}}(\bigoplus_{i\in E}\mathbf{M}_{i}) = \sup_{i\notin E}\mathrm{di}_{\mathbf{A}}(\mathbf{M}_{i})\) for every family \((\mathbf{M}_i)_{i\in E}\) of A-modules, A is noetherian (use Exercise 21, p.~170).
\end{enumerate}

\begin{enumerate}
\def\labelenumi{\arabic{enumi})}
\setcounter{enumi}{5}
\item
  Let I be the set of left ideals of A. Show that \(\mathrm{dh}(\mathbf{A}) = \sup \mathrm{dp}_{\mathbf{A}}(\mathbf{A}/\mathfrak{a})\) .
\item
  Let M be an A-module. We call the flat dimension of M, and denote it by \(dpl_{A}\) (M), the infimum in \(\overline{Z}\) of the lengths of the flat resolutions of M.
\end{enumerate}

\begin{enumerate}
\def\labelenumi{\alph{enumi})}
\tightlist
\item
  If n is an integer \(\geqslant 0\) show that the following conditions are equivalent:
\end{enumerate}

\(\alpha)\mathrm{dpl}_{\mathbf{A}}(\mathbf{M})\leqslant n.\)

\(\beta)\ \mathrm{Tor}_{r}^{\mathbf{A}}(\mathbf{N},\mathbf{M})=0\ \text{for every right A-module N and every integer }r>n.\)

\(\gamma)\ \mathrm{Tor}_{n+1}^{\mathbf{A}}(\mathbf{N},\mathbf{M})=0\) for every right A-module N.

\(\delta)\ \mathrm{Tor}_{n+1}^{\mathbf{A}}(\mathbf{N},\mathbf{M})=0\) for every monogenic right A-module N.

ε) For every exact sequence \(0 \rightarrow K \rightarrow P_{n-1} \rightarrow \cdots \rightarrow P_{0} \rightarrow M \rightarrow 0\) where the \(P_{k}\) are flat, K is flat.

\begin{enumerate}
\def\labelenumi{\alph{enumi})}
\setcounter{enumi}{1}
\item
  Show that one has \(\mathrm{dpl}_{\mathbf{A}}(\mathbf{M}) \leqslant \mathrm{dp}_{\mathbf{A}}(\mathbf{M})\) ; there is equality if A is noetherian and M is finitely generated.
\item
  Let \((\mathbf{N}_{\alpha})_{\alpha \in I}\) be an inductive system of A-modules and N its inductive limit. Prove the inequality
\end{enumerate}

\[
\mathrm{dpl} _ {A} (N) \leqslant \sup _ {\alpha \in I} \mathrm{dpl} _ {A} (N _ {\alpha}).
\]

\begin{enumerate}
\def\labelenumi{\arabic{enumi})}
\setcounter{enumi}{7}
\tightlist
\item
  We call the tor-dimension of A and denote by td(A) the supremum in \(\overline{\mathbf{Z}}\) of the set of integers \(n\) for which there exist a left A-module N and a right A-module M such that \(\operatorname{Tor}_{n}^{\mathbf{A}}(\mathbf{M},\mathbf{N})\neq0\) .
\end{enumerate}

\begin{enumerate}
\def\labelenumi{\alph{enumi})}
\tightlist
\item
  Let \(n\) be an integer \(\geqslant 0\) . Show that the following conditions are equivalent:
\end{enumerate}

\(\alpha)\ \mathrm{td}(\mathbf{A}) \leqslant n.\)

β) For every A-module M, we have \(\mathrm{dpl}_{\mathbf{A}}(\mathbf{M}) \leqslant n\) (Exercise 7).

\(\gamma)\) For every monogenic A-module M, we have \(\mathrm{dpl}_{\mathbf{A}}(\mathbf{M}) \leqslant n\) .

\begin{enumerate}
\def\labelenumi{\alph{enumi})}
\setcounter{enumi}{1}
\tightlist
\item
  Show that we have td(A) = td(A°) and td(A) ≤ dh(A); equality holds if A is left noetherian.
\end{enumerate}

¶ 9) a) Show that the following conditions are equivalent:

α) Every finitely generated (left) ideal of A is a projective A-module.

β) Every finitely generated submodule of a projective A-module is projective.

\(\gamma)\) We have td(A) \(\leqslant\) 1 (Exercise 8), and the ring A is coherent (X, p.~181, Exercise 11).

δ) For every set I, every submodule of A¹ is flat.

\begin{enumerate}
\def\labelenumi{\alph{enumi})}
\setcounter{enumi}{1}
\tightlist
\item
  Suppose that the ring A satisfies the equivalent conditions of a). Prove that every projective A-module is isomorphic to a direct sum of finitely generated ideals of A. (First treat the case of a finitely generated module; then treat by induction the case of a module admitting a countable family of generators, observing that every element of a projective A-module P is contained in a finitely generated direct summand of P. Deduce the general case from Kaplansky's theorem (II, p.~183, Exercise 2).)
\end{enumerate}

\begin{enumerate}
\def\labelenumi{\arabic{enumi})}
\setcounter{enumi}{9}
\tightlist
\item
  Assume \(\mathbf{A}\) is an integral domain; let \(\mathbf{K}\) be its field of fractions. We say that an ideal \(\mathfrak{a}\) of \(\mathbf{A}\) is \(invertible\) if there exist elements \(x_{1},\ldots ,x_{n}\) of \(\mathbf{K}\) such that \(x_{i}\mathfrak{a}\subset \mathbf{A}\) for \(1\leqslant i\leqslant n\) and \(\sum x_i\mathfrak{a} = \mathbf{A}\) .
\end{enumerate}

\begin{enumerate}
\def\labelenumi{\alph{enumi})}
\item
  Show that every invertible ideal is finitely generated.
\item
  Show that a nonzero ideal of A is a projective A-module if and only if it is invertible.
\item
  Let a be an invertible ideal of A, D a divisible A-module. Prove that \(\operatorname{Ext}_A^i(A/a, D) = 0\) for \(i \geqslant 1\) .
\end{enumerate}

\begin{enumerate}
\def\labelenumi{\arabic{enumi})}
\setcounter{enumi}{10}
\tightlist
\item
  \begin{enumerate}
  \def\labelenumii{\alph{enumii})}
  \tightlist
  \item
    Suppose the ring A is an integral domain. Show that the conditions \(\alpha\) to \(\delta\) of Exercise 9 are still equivalent to the following conditions:
  \end{enumerate}
\end{enumerate}

ε) Every finitely generated torsion-free A-module is projective.

ζ) Every torsion-free A-module is flat.

η) Every submodule of a flat A-module is flat.

(Use Exercise 13, p.~169.)

An integral domain satisfying these conditions is said to be Prüferian.

\begin{enumerate}
\def\labelenumi{\alph{enumi})}
\setcounter{enumi}{1}
\tightlist
\item
  Let B be an integral domain, the union of an increasing filtered family of Prüferian subrings. Show that B is Prüferian.
\end{enumerate}

* c) Show that the ring of algebraic integers is Prüferian. *

\begin{enumerate}
\def\labelenumi{\arabic{enumi})}
\setcounter{enumi}{11}
\tightlist
\item
  Suppose the ring A is an integral domain. Show that the following conditions are equivalent:
\end{enumerate}

\(\alpha)\ \mathrm{dh}(\mathbf{A}) \leqslant 1.\)

β) Every divisible A-module is injective.

γ) The ring A is Prüferian (Exercise 11) and noetherian.

(To show that α) implies β), use Exercise 10.)

We then say that A is a Dedekind ring.

\begin{enumerate}
\def\labelenumi{\arabic{enumi})}
\setcounter{enumi}{12}
\tightlist
\item
  Assume A is left and right noetherian. Show that the following conditions are equivalent:
\end{enumerate}

\(\alpha)\ \mathrm{dh}(\mathbf{A}) \leqslant 2.\)

β) For all integers \(p, q\) and every A-linear map \(u: \mathbf{A}^p \to \mathbf{A}^q\) , the kernel of \(u\) is a projective A-module.

γ) The dual of every finitely generated A-module is projective.

\begin{enumerate}
\def\labelenumi{\arabic{enumi})}
\setcounter{enumi}{13}
\tightlist
\item
  Let x be an element of A that is left-cancellable.
\end{enumerate}

\begin{enumerate}
\def\labelenumi{\alph{enumi})}
\item
  Show that the ideal \(x\mathbf{A}\) is two-sided if and only if there exists an endomorphism \(\sigma\) of the ring \(\mathbf{A}\) such that \(ax = x\sigma(a)\) for all \(a \in \mathbf{A}\) .
\item
  We now assume that the conditions of a) are satisfied. Let M be an A-module such that xM = 0. Prove the equality dp\_A(M) = dp\_{A/xA}(M) + 1 (one may use the spectral sequence of Exercise 6, d), p. 188).
\item
  We assume that \(x\) belongs to the radical of \(\mathbf{A}\) , and that \(\mathbf{A}\) is noetherian. Let \(\mathbf{N}\) be a finitely generated \(\mathbf{A}\) -module, such that the homothety of ratio \(x\) is injective in \(\mathbf{N}\) . Show that we have \(\mathrm{dp}_{\mathbf{A}}(\mathbf{N}) = \mathrm{dp}_{\mathbf{A}/\mathbf{A}x}(\mathbf{N}/x\mathbf{N})\) . d) Under the conditions of \(c\) ), show that one has \(\mathrm{dh}(\mathbf{A}) = \mathrm{dh}(\mathbf{A}/\mathbf{A}x) + 1\) .
\end{enumerate}

\begin{enumerate}
\def\labelenumi{\arabic{enumi})}
\setcounter{enumi}{14}
\tightlist
\item
  Let σ be an automorphism of the ring A.
\end{enumerate}

\begin{enumerate}
\def\labelenumi{\alph{enumi})}
\tightlist
\item
  We denote by \(\mathbf{A}_{\sigma}[\mathbf{X}]\) the ring defined in VIII, § 1, no. 3. Show that
\end{enumerate}

\[
\mathrm{dh} (\mathbf {A} _ {\sigma} [ \mathbf {X} ]) = \mathrm{dh} (\mathbf {A}) + 1.
\]

(The inequality dh(A \(_{\sigma}\) {[}X{]}) \(\geqslant\) dh(A) + 1 follows from Exercise 14, \(b\) ); for the opposite inequality, be inspired by the proof of X, p.~142, lemma 3, \(a\) .)

\begin{enumerate}
\def\labelenumi{\alph{enumi})}
\setcounter{enumi}{1}
\tightlist
\item
  We denote by \(\mathbf{A}_{\sigma}[[\mathbf{X}]]\) the ring defined in IV, § 4, Exercise 8. Show that if A is noetherian, then we have \(\mathrm{dh}(\mathbf{A}_{\sigma}[[\mathbf{X}]))=\mathrm{dh}(\mathbf{A})+1\) (use Exercise 14).
\end{enumerate}

\begin{enumerate}
\def\labelenumi{\arabic{enumi})}
\setcounter{enumi}{15}
\tightlist
\item
  Let A be a local ring that is noetherian on the right and on the left, and let m be its maximal ideal; we denote by \(k_{s}\) (resp. \(k_{d}\) ) the left (resp. right) A-module A/m. Let M be a Tor \(_{n+1}^{A}\) finitely generated left module and let n be an integer. a) Show that one has dp \(_{A}\) (M) \(\leqslant\) n if and only if Tor \(_{n+1}^{A}\) ( \(k_{d}\) , M) = 0 (use Exercise 4, p.~185).
\end{enumerate}

\begin{enumerate}
\def\labelenumi{\alph{enumi})}
\setcounter{enumi}{1}
\tightlist
\item
  Show that one has \(\mathrm{dh(A)} = \mathrm{dp}_{\mathbf{A}}(k_s)\) ; for \(\mathrm{dh(A)} \leqslant n\) it is necessary and sufficient that one has
\end{enumerate}

\[
\operatorname{Tor} _ {n + 1} ^ {\mathbf {A}} \left(k _ {d}, k _ {s}\right) = 0.
\]

\begin{enumerate}
\def\labelenumi{\alph{enumi})}
\setcounter{enumi}{2}
\tightlist
\item
  Let \(x\) be an element of the center of \(\mathbf{A}\) belonging to \(m\) , such that the homothety with ratio \(x\) in \(M\) is injective. Show that we have \(\mathrm{dp}_{\mathrm{A}}(\mathrm{M}/x\mathrm{M}) = \mathrm{dp}_{\mathrm{A}}(\mathrm{M}) + 1\) (cf. also Exercise 14).
\end{enumerate}

\begin{enumerate}
\def\labelenumi{\arabic{enumi})}
\setcounter{enumi}{16}
\tightlist
\item
  Let M be an A-module, I a well-ordered set, \((\mathbf{M}_{\alpha})_{\alpha \in I}\) an increasing family of submodules of M, such that \(M = \bigcup M_{\alpha}\) . We set \(M_{\alpha}^{\prime} = \bigcup M_{\beta}\) for \(\alpha \in I\) .
\end{enumerate}

\begin{enumerate}
\def\labelenumi{\alph{enumi})}
\item
  Prove the inequality dpA(M) ≤ sup dpA(Mα/Mα').
\item
  Suppose there exists an integer \(n\) such that \(\mathrm{dp}_{\mathbf{A}}(\mathbf{M}_{\alpha}') \leqslant n\) for all \(\alpha \in I\) . Show that we have
\end{enumerate}

\[
\mathrm{dp} _ {\mathbf {A}} (\mathbf {M}) \leqslant n + 1.
\]

\begin{enumerate}
\def\labelenumi{\arabic{enumi})}
\setcounter{enumi}{17}
\tightlist
\item
  \begin{enumerate}
  \def\labelenumii{\alph{enumii})}
  \tightlist
  \item
    Let \((\mathbf{M}_{\alpha})_{\alpha \in I}\) be an inductive system of A-modules relative to a countable set I, M its limit, \(n\) be an integer \(\geqslant 0\) . Show that if \(\mathrm{dp}_{\mathbf{A}}(\mathbf{M}_{\alpha}) \leqslant n\) for all \(\alpha \in I\) , we have \(\mathrm{dp}_{\mathbf{A}}(\mathbf{M}) \leqslant n + 1\) (reduce to the case \(n = 0\) and \(I = N\) , then write an exact sequence \(0 \to \bigoplus_{n} \mathbf{M}_{n} \to \bigoplus_{n} \mathbf{M}_{n} \to \mathbf{M} \to 0\) ).
  \end{enumerate}
\end{enumerate}

\begin{enumerate}
\def\labelenumi{\alph{enumi})}
\setcounter{enumi}{1}
\item
  Let P be a flat A-module admitting a presentation \(\mathbf{A}^{(\mathbb{S})}\to\mathbf{A}^{(\mathbb{T})}\to\mathbf{P}\to0\) , where the set S is countable. Deduce from a) that one has dp \(_{\mathbb{A}}(\mathbb{P})\leqslant1\) .
\item
  Assume that every left ideal of the ring A is generated by a countable set of elements. Prove that dpA(M) ≤ dplA(M) + 1 for every A-module M generated by a countable family of elements. Deduce the inequalities dh(A) ≤ td(A) + 1 and \textbar dh(A) - dh(A°)\textbar{} ≤ 1. (Use Exercise 12, p.~181, Exercise 8, p.~202, as well as b).)
\end{enumerate}

¶ 19) a) Let n be an integer ≥ 0, let (Mα)α ∈ I be an inductive system of A-modules, and let M be its inductive limit. Assume that Card (I) ≤ Nn (E, III, p.~87, exercise 10) and that there exists an integer r such that dpA(Mα) ≤ r for every α ∈ I. Show that dpA(M) ≤ r + n + 1 (argue by induction on n, using exercise 18, a) and exercise 17).

\begin{enumerate}
\def\labelenumi{\alph{enumi})}
\setcounter{enumi}{1}
\tightlist
\item
  We assume that every left ideal of A is generated by a family of elements of cardinal \(\leqslant \aleph_{n}\) . Prove that one has dp \(_{A}\) (M) \(\leqslant\) dpl \(_{A}\) (M) + n + 1 for every A-module M generated by a family of elements of cardinality \(\leqslant \aleph_{n}\) . Deduce from this the inequalities
\end{enumerate}

\[
\mathrm{dh} (\mathbf {A}) \leqslant \mathrm{td} (\mathbf {A}) + n + 1 \quad \text { and } \quad | \mathrm{dh} (\mathbf {A}) - \mathrm{dh} (\dot {\mathbf {A}} ^ {\circ}) | \leqslant n + 1.
\]

¶ 20) Let R be a principal ideal ring, K its field of fractions. We denote by B the subring of M₂(K) consisting of the matrices \(\begin{pmatrix} a & 0 \\ x & y \end{pmatrix}\) with \(a \in \mathbb{R}, x, y \in \mathbb{K}\) .

\begin{enumerate}
\def\labelenumi{\alph{enumi})}
\item
  Show that B is left noetherian but not right noetherian.
\item
  Show that \(\mathrm{dh(B)} = 1\) and \(\mathrm{dh(B^{\circ})} = 2\) .
\end{enumerate}

\begin{enumerate}
\def\labelenumi{\arabic{enumi})}
\setcounter{enumi}{20}
\tightlist
\item
  Let r be the radical of A. Show that the following conditions are equivalent:
\end{enumerate}

α) A/r is semisimple, and for every sequence \((a_{n})_{n\geq1}\) of elements of r, there exists an integer n such that \(a_{1}\ldots a_{n}=0\) .

β) Every A-module admits a projective cover.

\(\gamma)\) For every A-module M, we have \(\mathrm{dpl}_{\mathbf{A}}(\mathbf{M}) = \mathrm{dp}_{\mathbf{A}}(\mathbf{M})(\mathbf{X}, p. 202, \text{exercise } 7)\) .

δ) For every inductive system \((\mathbf{M}_{\alpha}, f_{\beta\alpha})_{\alpha \in I}\) of A-modules, we have:

\[
\mathrm{dp} _ {\mathrm{A}} (\varinjlim \mathrm{M} _ {\alpha}) \leqslant \sup _ {\alpha \in I} \mathrm{dp} _ {\mathrm{A}} (\mathrm{M} _ {\alpha}).
\]

ε) Every flat A-module is projective.

φ) Every descending sequence of monogenic right ideals of A is stationary.

(For the equivalence of \(\alpha\) ) and \(\beta\) ), cf.~VIII, § 8, exercise 26. To show that \(\beta\) ) implies \(\gamma\) ), note that M admits a minimal projective resolution P and that \(P_n/rP_n\) is isomorphic to Tor \(_n^A\) (A/r, M). To prove \(\varepsilon) \Rightarrow \varphi\) ), let \((l_n)_{n \geq 1}\) be a decreasing sequence of monogenic right ideals, such that

\[
\mathrm{I} _ {n} = a _ {1} \dots a _ {n} \mathrm{A}, \quad \text { with } \quad a _ {i} \in \mathrm{A}.
\]

In the free module \(\mathbf{A}^{(\mathbf{N})}\) , consider the submodules \(\mathbf{N}_p\) generated by \(e_1 - a_1 e_2, \ldots, e_p - a_p e_{p+1}\) . Deduce from \(\varepsilon\) ) that \(\mathbf{N} = \bigcup \mathbf{N}_p\) is a direct summand in \(\mathbf{A}^{(\mathbf{N})}\) ; conclude from this that \(\mathrm{l}_{n+1} = \mathrm{l}_n\) for \(n\) large enough.

For the implication \(\varphi) \Rightarrow \alpha)\) , use VIII, § 9, exercise 12.)

\begin{enumerate}
\def\labelenumi{\arabic{enumi})}
\setcounter{enumi}{21}
\tightlist
\item
  Show that the following conditions are equivalent:
\end{enumerate}

(α) Every finitely presented A-module of finite projective dimension is projective.

(β) Every injective A-homomorphism \(A^{p} \rightarrow A^{q}\) admits a retraction.

(γ) For every finitely generated right ideal a distinct from A, there exists a nonzero element x of A such that xa = 0.

(δ) Every simple right A-module of finite presentation is isomorphic to a (minimal) right ideal of A.

(Observe that condition \(\gamma\) ) is equivalent to condition \(\beta\) in the case \(p = 1\) .)

¶ 23) Show that the following conditions are equivalent:

(α) Every A-module of finite projective dimension is projective.

(β) The ring A satisfies the equivalent conditions of Exercise 21, and for every finitely generated right ideal a distinct from A, there exists an element \(x \neq 0\) of A such that \(xa = 0\) .

(γ) The ring A satisfies the equivalent conditions of Exercise 21, and every simple A-module is a quotient of an injective A-module.

(Show that \(\alpha\) ) implies that every descending sequence of monogenic right ideals of A is stationary, by adapting the proof of Exercise 21. Then deduce from Exercise 22 the equivalence of \(\alpha\) ) and \(\beta\) ). To see that \(\alpha\) ) implies \(\gamma\) ), consider a projective cover P of the simple module M, an injective module I containing P, and a projective cover Q of I; define an A-homomorphism from Q to P, and deduce from it a surjection from I to M. To show that \(\gamma\) ) implies \(\alpha\) ), show that an A-module M such that dp \(_{A}\) (M) \(\leqslant\) 1 is projective by considering a projective cover \(p: P \to M\) and a surjection from Ker \(p\) onto a simple module.)

\begin{enumerate}
\def\labelenumi{\arabic{enumi})}
\setcounter{enumi}{23}
\tightlist
\item
  We denote by dhf(A) (resp. dif(A), resp. tdf(A)) the supremum of the projective (resp. injective, resp. flat) dimensions of A-modules for which this dimension is finite. If A has finite homological dimension, then dhf(A) = dif(A) = dh(A) and tdf(A) = td(A). If td(A) \textless{} ∞, we have tdf(A) = td(A). a) Show that tdf(A) ≤ dif(A°); equality holds if A is left noetherian (use X, p.~189, exercise 3).
\end{enumerate}

\begin{enumerate}
\def\labelenumi{\alph{enumi})}
\setcounter{enumi}{1}
\item
  Assume A is left noetherian. Show that we have dhf(A) ≤ di\_A(A\_s). If di\_A(A\_s) and dif(A) are finite, show that they are equal (use Exercise 2, p.~202).
\item
  Let G be a finite group not reduced to the identity element, and let A be the ring \(\mathbf{Z}^{(G)}\) . Show that we have
\end{enumerate}

\[
\mathrm{dh} (\mathbf {A}) = + \infty , \mathrm{dhf} (\mathbf {A}) = \mathrm{dif} (\mathbf {A}) = \mathrm{tdf} (\mathbf {A}) = 1, \mathrm{di} _ {\mathbf {A}} (\mathbf {A} _ {s}) = + \infty
\]

(cf. exercise 20, p. 193).

\begin{enumerate}
\def\labelenumi{\alph{enumi})}
\setcounter{enumi}{3}
\tightlist
\item
  Assume the ring A is self-injective (cf. X, p. 171, exercise 26), not semisimple. Show that every non-projective A-module has infinite projective dimension and infinite injective dimension. Deduce that one has dh(A) = +∞, dhf(A) = dif(A) = tdf(A) = di\_A(A\_s) = 0.
\end{enumerate}

\begin{enumerate}
\def\labelenumi{\arabic{enumi})}
\setcounter{enumi}{24}
\tightlist
\item
  We denote by \(\mathbf{A}^e\) the \(k\) -algebra \(\mathbf{A} \otimes_k \mathbf{A}^*\) and we consider \(\mathbf{A}\) as a left \(\mathbf{A}^e\) -module. We set
\end{enumerate}

\[
\mathrm{da} _ {k} (\mathrm{A}) = \mathrm{dp} _ {\mathrm{A} ^ {e}} (\mathrm{A}).
\]

If n is an integer \(\geqslant0\) , the inequality \(\mathrm{da}_{k}(\mathrm{A})\leqslant n\) is therefore equivalent to the vanishing of \(\mathrm{H}^{n+1}(\mathrm{A},\mathrm{M})\) for every (A, A)-bimodule M (X, p.~195, exercise 26). One has \(\mathrm{da}_{k}(\mathrm{A})=\mathrm{da}_{k}(\mathrm{A}^{\circ})\) .

\begin{enumerate}
\def\labelenumi{\alph{enumi})}
\item
  Show that one has \(\mathrm{da}_k(\mathbf{M}_n(k)) = 0\) and \(\mathrm{da}_k(k[\mathbf{X}_1, \ldots, \mathbf{X}_n]) = n\) for all \(n \geqslant 0\) .
\item
  If \(k\) is a field, one has \(\mathrm{da}_{k}(\mathbf{A}) = 0\) if and only if the \(k\) -algebra \(\mathbf{A}\) is absolutely semisimple (cf.~VIII, § 11, n° 3, th. 2).
\item
  Let M be a free \(k\) -module. Prove that one has \(\mathrm{da}_{k}(\mathbf{T}_{k}(\mathbf{M})) = 1\) (If \((e_{i})_{i \in I}\) is a basis of M, prove that the elements \((e_{i} \otimes 1 - 1 \otimes e_{i})\) form a basis of the left ideal of \(\mathbf{A}^{e}\) kernel of multiplication \(\mathbf{A} \otimes_{k} \mathbf{A}^{-} \to \mathbf{A}\) , with \(\mathbf{A} = \mathbf{T}_{k}(\mathbf{M})\) ).
\item
  We now assume that the \(k\) -module A is projective. Let \(\rho: k \to k'\) be a homomorphism of commutative rings; show that one has \(\mathrm{da}_k(\mathbf{A} \otimes_k k') \leqslant \mathrm{da}_k(\mathbf{A})\) . If \(\rho\) is injective and if \(\rho(k)\) is a direct summand of the \(k\) -module \(k'\) , one has equality (use X, p.~195, exercise 26, c)).
\item
  Let B be a second k-algebra (associative and unital), projective over k. Show that
\end{enumerate}

\[
\mathrm{da} _ {k} (\mathbf {A} \otimes_ {k} \mathbf {B}) \leqslant \mathrm{da} _ {k} (\mathbf {A}) + \mathrm{da} _ {k} (\mathbf {B}).
\]

If moreover k is a field and A and B are finite-dimensional over k, then equality holds (use X, p.~195, exercise 26, d)).

\begin{enumerate}
\def\labelenumi{\alph{enumi})}
\setcounter{enumi}{5}
\tightlist
\item
  Assume that k is a field. Prove the inequalities \(\mathrm{dh}(\mathbf{A}) \leqslant \mathrm{da}_{k}(\mathbf{A})\) and \(\mathrm{dh}(\mathbf{A}^{\circ}) \leqslant \mathrm{da}_{k}(\mathbf{A})\) (cf. X, p. 196, exercise 26, e)).
\end{enumerate}

\begin{enumerate}
\def\labelenumi{\arabic{enumi})}
\setcounter{enumi}{25}
\tightlist
\item
  Let A be an augmented \(k\) -algebra (X, p.~196, exercise 27); one assumes that there exists a homomorphism \(\rho: \mathrm{A} \to \mathrm{A}^e\) which satisfies the conditions of X, p.~196, exercise 27, c). Show that \(\mathrm{da}_k(\mathbf{A}) = \mathrm{dp}_\mathbf{A}(k)\) ; if moreover \(k\) is a field, one has \(\mathrm{da}_k(\mathbf{A}) = \mathrm{dp}_\mathbf{A}(k) = \mathrm{dh}(\mathbf{A}) = \mathrm{dh}(\mathbf{A}^\circ)\) . (Use Exercise 27, b) and c), p.~196, and Exercise 26, e), p.~196.) In particular, if V is a nonzero \(k\) -vector space, we have \(\mathrm{dh}\,\mathbf{T}_k(\mathbf{V}) = 1\) .
\end{enumerate}

* 27) We assume that \(k\) is a field. Let \(g\) be a Lie algebra over \(k\) of dimension \(n\) , U its enveloping algebra. Show that \(dhU = n\) . (Use Exercise 26, Exercise 27, d), p.~196, and Exercise 28, c), p.~197.) *

¶ 28) Let G be a group. One calls the cohomological dimension of G, denoted dc(G), the projective dimension of the \(\mathbf{Z}^{(G)}\) -module Z. If \(n\) is an integer, we thus have dc(G) \(\leqslant n\) if and only if H \(^{q}\) (G, M)=0 for every \(\mathbf{Z}^{(G)}\) -module M and every \(q > n\) .

\begin{enumerate}
\def\labelenumi{\alph{enumi})}
\item
  Prove that one has \(\mathrm{da}_{\mathbf{Z}}(\mathbf{Z}^{(G)}) = \mathrm{dc}(G)\) (use Exercise 26 and Exercise 27, d), p.~196).
\item
  If G is a free group not reduced to the identity element, then dc(G) = 1.
\item
  Let H be a subgroup of G; prove the inequality dc(H) ≤ dc(G). If dc(G) \textless{} ∞ and if H has finite index in G, show that dc(H) = dc(G) (show that for q = dc(G), the homomorphism
\end{enumerate}

\[
t ^ {q}: \mathrm{H} ^ {q} (\mathrm{H}, \mathrm{M}) \rightarrow \mathrm{H} ^ {q} (\mathrm{G}, \mathrm{M})
\]

of Exercise 14, p.~191, is surjective).

\begin{enumerate}
\def\labelenumi{\alph{enumi})}
\setcounter{enumi}{3}
\item
  Show that the cohomological dimension of a finite group not reduced to the identity element is infinite. Deduce that if dc(G) \textless{} ∞, every element of G other than the identity element has infinite order (one then says that G is torsion-free).
\item
  If G is torsion-free and if H is a subgroup of finite index in G, show that dc(H) = dc(G). (If (P, p) is a projective resolution of \(\mathbf{Z}^{(H)}\) -module Z and if (G : H) = n, define on the graded \(\mathbf{Z}^{(H)}\) -module \(\mathbf{T}_{\mathbf{Z}}^{n}(\mathbf{P})\) a structure of \(\mathbf{Z}^{(G)}\) -complex so that \((\mathbf{T}_{\mathbf{Z}}^{n}(\mathbf{P}), p^{\otimes n})\) is a projective resolution of the \(\mathbf{Z}^{(G)}\) -module Z.)
\item
  If H is a normal subgroup of G, show that we have
\end{enumerate}

\[
\mathrm{dc} (\mathbf {G}) \leqslant \mathrm{dc} (\mathbf {H}) + \mathrm{dc} (\mathbf {G} / \mathbf {H}).
\]

\exercisesection{Koszul Complexes}

\begin{enumerate}
\def\labelenumi{\arabic{enumi})}
\tightlist
\item
  Let A be a ring, L an A-module, and K the complex \(\mathbf{S}(\mathbf{L})\otimes_{\mathbf{A}}\symbf{\Lambda}(\mathbf{L})\) (X, p.~151, example 1).
\end{enumerate}

\begin{enumerate}
\def\labelenumi{\alph{enumi})}
\tightlist
\item
  Show that there exists an A-linear endomorphism s of K, graded of degree zero, such that
\end{enumerate}

\[
d s (x) + s d (x) = (p + q) x
\]

for all integers \(p, q \geqslant 0\) and every \(x \in \mathbf{S}^p L \otimes_{\mathbf{A}} \Lambda^q L\) .

\begin{enumerate}
\def\labelenumi{\alph{enumi})}
\setcounter{enumi}{1}
\tightlist
\item
  Assume that A is a Q-algebra. Show that K defines a resolution of the A-module A.
\end{enumerate}

\begin{enumerate}
\def\labelenumi{\arabic{enumi})}
\setcounter{enumi}{1}
\tightlist
\item
  Let A be a ring, \(u: \mathbf{L} \to \mathbf{M}\) a homomorphism of A-modules.
\end{enumerate}

\begin{enumerate}
\def\labelenumi{\alph{enumi})}
\item
  We denote by \(\varepsilon\) the commutation factor on \(\mathbf{Z}^2\) defined by \(\varepsilon(\alpha, \beta) = \alpha_2 \beta_2\) for \(\alpha, \beta \in \mathbf{Z}^2\) , \(\alpha = (\alpha_1, \alpha_2)\) , \(\beta = (\beta_1, \beta_2)\) . Show that there exists a unique (A, \(\varepsilon\) )-derivation (III, p.~118) \(d\) of the bigraded algebra \(\mathbf{S}(M) \otimes_A \symbf{\Lambda}(L)\) of degree \((1, -1)\) , such that one has \(d(1 \otimes x) = u(x) \otimes 1\) for \(x \in L\) . Show that \(d \circ d = 0\) , so that \(\mathbf{S}(M) \otimes_A \symbf{\Lambda}(L)\) , equipped with the grading induced by that of \(\symbf{\Lambda}(L)\) and with the differential \(d\) , defines a complex \(\mathbf{K}(u)\) .
\item
  If M = L and u = \(ld_{L}\) , show that \(\mathbf{K}(u)\) is equal to the complex defined in no. 3, p.~151.
\item
  Let \(\mathbf{B} = \mathbf{S}_{\mathbf{A}}(\mathbf{M})\) , and let \(\overline{u}: \mathbf{B} \otimes_{\mathbf{A}} \mathbf{L} \to \mathbf{B}\) be the B-homomorphism deduced from \(u\) . Show that \(\mathbf{K}(u)\) is canonically identified with the complex \(\mathbf{K}^{\mathbf{B}}(\overline{u})\) .
\item
  Let \(u': L' \to M'\) be a second A-homomorphism. Show that the complexes \(K(u \oplus u')\) and \(K(u) \otimes_{A} K(u')\) are isomorphic.
\item
  Let \(f: \mathbf{L} \to \mathbf{L}'\) and \(g: \mathbf{M} \to \mathbf{M}'\) be two A-homomorphisms such that \(g \circ u = u' \circ f\) . Show that the homomorphism \(\mathbf{S}(g) \otimes \symbf{\Lambda}(f)\) is a morphism of complexes from \(\mathbf{K}(u)\) to \(\mathbf{K}(u')\) . In particular, if \(v: \mathbf{M} \to \mathbf{P}\) is an A-homomorphism such that \(v \circ u = 0\) , we obtain a morphism of complexes \(h: \mathbf{K}(u) \to \mathbf{S}_{\mathbf{A}}(\mathbf{P})\) .
\item
  We assume that the sequence \(0 \rightarrow L \xrightarrow{u} M \xrightarrow{v} P \rightarrow 0\) is an exact sequence of flat A-modules. Show that h is an isomorphism (reduce to the case where L, M, P are finitely generated free, and use d)). Deduce from this for every \(n \geqslant 1\) an exact sequence
\end{enumerate}

\[
0 \rightarrow \Lambda^ {n} L \rightarrow M \otimes_ {A} \Lambda^ {n - 1} L \rightarrow \dots \rightarrow S ^ {n - 1} M \otimes_ {A} L \rightarrow S ^ {n} M \rightarrow S ^ {n} P \rightarrow 0.
\]

\begin{enumerate}
\def\labelenumi{\arabic{enumi})}
\setcounter{enumi}{2}
\tightlist
\item
  Let \(n\) be an integer, \(k\) a ring, \(A = k[X_1, \ldots, X_n]\) . Show that for every A-module M there exists an isomorphism of graded A-modules
\end{enumerate}

\[
\delta_ {\mathbf {M}}: \operatorname{Tor} ^ {\mathbf {A}} (k, \mathbf {M}) \rightarrow \operatorname{Ext} _ {\mathbf {A}} (k, \mathbf {M}) (- n)
\]

such that for every homomorphism of A-modules \(u: M \to N\) , one has Ext \((1_k, u) \circ \delta_M = \delta_N \circ \text{Tor } (1_k, u)\) . 4) Let A be a ring, L an A-module, \(u: L \to A\) a linear form.

\begin{enumerate}
\def\labelenumi{\alph{enumi})}
\item
  Show that the complex \(\mathbf{K}(u)\) , endowed with the algebra structure of \(\symbf{\Lambda}(\mathrm{L})\) , is a graded differential A-algebra (X, p.~183, exercise 18).
\item
  Let B be a differential graded A-algebra such that \(B_{0} = A\) and let \(f : L \to B_{1}\) be an A-homomorphism such that \(d_{1} \circ f = u\) . Show that there exists a unique morphism of graded differential A-algebras \(F : \mathbf{K}(u) \to \mathbf{B}\) such that \(F_{0} = 1_{A}\) and \(F_{1} = f\) .
\end{enumerate}

\begin{enumerate}
\def\labelenumi{\arabic{enumi})}
\setcounter{enumi}{4}
\item
  Let A be a ring, M an A-module, \(x = (x_1, \ldots, x_n)\) a sequence of elements of A, \(k = (k_1, \ldots, k_n) \in \mathbf{N}^n\) , \(x^k = (x_1^{k_1}, \ldots, x_{n}^{k_n})\) . Prove that the sequence \(x^k\) is M-regular if and only if the sequence \(x\) is; in that case the A-module \(\mathrm{M}/(x^k)\) admits a composition series whose quotients are isomorphic to \(\mathrm{M}/(x)\) . (Argue by induction on \(n\) .)
\item
  Let A be a ring, \(x = (x_1, \ldots, x_n)\) a sequence of elements of A, \(x\) the ideal \(\sum_{i} x_i\) A. Let
\end{enumerate}

\[
0 \rightarrow \mathbf {M} ^ {\prime} \stackrel {i} {\rightarrow} \mathbf {M} \rightarrow \mathbf {M} ^ {\prime \prime} \rightarrow 0
\]

an exact sequence of A-modules. Suppose that the sequence x is completely secant for M and that the A-modules \(\mathrm{M}''/(x_{1}\mathrm{M}''+\cdots+x_{i-1}\mathrm{M}'')\) are separated for the x-adic topology ( \(1\leqslant i\leqslant n\) ). Show that the following conditions are equivalent:

α) The sequence x is completely secant for M''.

β) The homomorphism \(\bigoplus_{r} (\mathfrak{x}^{r} \, \mathrm{M}^{\prime}/\mathfrak{x}^{r+1} \, \mathrm{M}') \to \bigoplus_{r} (\mathfrak{x}^{r} \, \mathrm{M}/\mathfrak{x}^{r+1} \, \mathrm{M})\) induced by \(i\) is injective.

\(\gamma)\) The homomorphism \(\mathbf{M}' / \mathfrak{x}\mathbf{M}' \to \mathbf{M} / \mathfrak{x}\mathbf{M}\) induced by \(i\) is injective.

\begin{enumerate}
\def\labelenumi{\arabic{enumi})}
\setcounter{enumi}{6}
\tightlist
\item
  Let A be a ring, M an A-module, \(x = (x_1, \ldots, x_n)\) a sequence of elements of A.
\end{enumerate}

\[
(X, p p. 1 7 5 - 1 7 7,
\]

\[
\mathrm{E} _ {p q} ^ {2} = \operatorname{Tor} _ {p} ^ {\mathrm{A}} \left(\mathrm{H} _ {q} (\boldsymbol {x}, \mathrm{A}) \right.
\]

\begin{enumerate}
\def\labelenumi{\alph{enumi})}
\setcounter{enumi}{1}
\tightlist
\item
  Let \(p\) be an integer \(\leqslant n\) ; set \(x' = (x_1, \ldots, x_p)\) and \(x'' = (x_{p+1}, \ldots, x_n)\) . Show that there exists a spectral sequence E such that 'E \(_{pa}^{2}\) = H \(_{p}\) x'`, H \(_{a}\) (x', M)), converging to H. (x, M).
\end{enumerate}

\begin{enumerate}
\def\labelenumi{\arabic{enumi})}
\setcounter{enumi}{7}
\tightlist
\item
  Let A be a ring, M an A-module, \(\boldsymbol{x}' = (x_1, \ldots, x_p)\) and \(\boldsymbol{x}'' = (x_{p+1}, \ldots, x_n)\) two sequences of elements of A; we denote by \(\boldsymbol{x}\) the sequence \((x_1, \ldots, x_n)\) , and by \(\boldsymbol{M}'\) the A-module \(\mathrm{M}/(\boldsymbol{x}')\) M.
\end{enumerate}

\begin{enumerate}
\def\labelenumi{\alph{enumi})}
\item
  If \(x'\) is completely secant for M and \(x''\) completely secant for M', show that x is completely secant for M.
\item
  If the sequences \(x\) and \(x'\) are completely secant for M, show that the sequence \(x''\) is completely secant for M'.
\item
  Give an example of an A-regular sequence \((x, y)\) contained in the radical of A, such that x is completely secant for A/yA, but y is not completely secant for A.
\end{enumerate}

\begin{enumerate}
\def\labelenumi{\arabic{enumi})}
\setcounter{enumi}{8}
\tightlist
\item
  Let A be a ring, M an A-module, and k an integer \(\geqslant\) 1. One says that a family x of elements of A is k-secant for M if one has \(H_{i}(x, M) = 0\) for \(1 \leqslant i \leqslant k\) . A family x is therefore completely secant if and only if it is k-secant for every \(k \geqslant 1\) .
\end{enumerate}

\begin{enumerate}
\def\labelenumi{\alph{enumi})}
\item
  Let \(0 \to M' \to M \to M'' \to 0\) be an exact sequence of A-modules. Show that if the family \(x\) is \(k\) -secant for \(M'\) and \(M''\) , it is so for \(M\) .
\item
  If x is k-secant for M and if N is a flat A-module, then x is k-secant for M \(\otimes_{A} N\) .
\item
  Let \(a_{1},\ldots,a_{n}\) be integers \(\geqslant 1\) . Show that a sequence \((x_{1},\ldots,x_{n})\) is \(k\) -secant for M if and only if the same holds for the sequence \((x_{1}^{a_{1}},\ldots,x_{n}^{a_{n}})\) .
\item
  For every pair of integers k, n with \(1 \leqslant k < n\) , give an example of a ring A, an A-module M, and a sequence \((x_{1}, \ldots, x_{n})\) of elements of A that is k-secant but not \((k + 1)\) -secant for M (use X, p.~159, remark 4).
\end{enumerate}

\begin{enumerate}
\def\labelenumi{\arabic{enumi})}
\setcounter{enumi}{9}
\tightlist
\item
  Let A be a ring, M an A-module, \(\boldsymbol{x} = (x_{1}, \ldots, x_{n})\) a sequence of elements of A, p an integer \(\leqslant n\) . We write \(\boldsymbol{x}' = (x_{1}, \ldots, x_{p})\) , \(\boldsymbol{x}'' = (x_{p+1}, \ldots, x_{n})\) and \(\mathrm{M}' = \mathrm{M}/(\boldsymbol{x}')\) M.
\end{enumerate}

\begin{enumerate}
\def\labelenumi{\alph{enumi})}
\item
  Let \(k, k'\) be two integers such that \(k \leqslant k'\) . Suppose that the sequence \(x'\) is \(k'\) -secant for M (exercise 9). Show that the sequence \(x\) is \(k\) -secant for M if and only if the sequence \(x''\) is \(k\) -secant for M'.
\item
  Show that if the sequence x is 1-secant for M, then the sequence \(x''\) is 1-secant for \(M'\) .
\end{enumerate}

¶ 11) Let A be a noetherian local ring, m its maximal ideal, k = A/m. Let M be a nonzero finitely generated A-module. The depth of M, denoted prof (M), is defined as the supremum of the integers n for which there exists a sequence \((x_{1}, \ldots, x_{n})\) of elements of m that is completely secant for M.

\begin{enumerate}
\def\labelenumi{\alph{enumi})}
\item
  Show that prof (M) = 0 if and only if m is associated to M (X, p.~172, exercise 28), in other words, if M contains a submodule isomorphic to k (observe that an ideal contained in a union of prime ideals is necessarily contained in one of them).
\item
  Show that prof(M) is the infimum of the integers n for which Ext \(_{A}^{n}\) (k, M) is nonzero (argue by induction on n). If prof(M) = n, then every sequence (x \(_{1}\) , \ldots, x \(_{p}\) ) of elements of m that is completely secant for M (p ≤ n) can be extended to a sequence (x \(_{1}\) , \ldots, x \(_{n}\) ) in m completely secant for M. c) If (x \(_{1}\) , \ldots, x \(_{k}\) ) is a sequence of elements of m completely secant for M, show that
\end{enumerate}

\[
\operatorname{prof} \left(\mathbf {M} / (x _ {1} \mathbf {M} + \dots + x _ {k} \mathbf {M})\right) = \operatorname{prof} (\mathbf {M}) - k.
\]

\begin{enumerate}
\def\labelenumi{\alph{enumi})}
\setcounter{enumi}{3}
\item
  Assume that M has finite projective dimension. Prove the equality dpA(M) + prof(M) = prof(A) (argue by induction on dpA(M)).
\item
  Show that prof(A) is the supremum of the projective dimensions of A-modules that are finitely generated and of finite projective dimension.
\end{enumerate}

¶ 12) Let A be a noetherian local ring, m its maximal ideal.

\begin{enumerate}
\def\labelenumi{\alph{enumi})}
\tightlist
\item
  Show that if \(0 < \mathrm{dh}(\mathbf{A}) < \infty\) there exists a non-zero-divisor element in \(\mathfrak{m} - \mathfrak{m}^2\) (show as in Exercise 11, a) that \(\mathfrak{m}\) would otherwise be associated with \(\mathbf{A}\) whence an exact sequence
\end{enumerate}

\[
0 \to k \to \mathbf {A} \to \mathbf {N} \to 0;
\]

lead to a contradiction by using Exercise 16, p.~203).

\begin{enumerate}
\def\labelenumi{\alph{enumi})}
\setcounter{enumi}{1}
\tightlist
\item
  For any integer n, show that dh(A) = n if and only if there exists a sequence \((x_{1}, \ldots, x_{n})\) completely secant for A, generating m (argue by induction on n, using a) and exercise 14, p.~203).
\end{enumerate}

A local ring is said to be regular if it is noetherian and has finite homological dimension.

\begin{enumerate}
\def\labelenumi{\arabic{enumi})}
\setcounter{enumi}{12}
\tightlist
\item
  Let A be a regular local ring (Exercise 14) of homological dimension n; denote its maximal ideal by m and let k = A/m.
\end{enumerate}

\begin{enumerate}
\def\labelenumi{\alph{enumi})}
\item
  Show that one has \(\dim_k(m/m^2) = n\) .
\item
  Show that the \(k\) -algebra \(\bigoplus_{r\geqslant 0}(m^{r} / m^{r + 1})\) is isomorphic to a polynomial algebra in \(n\) indeterminates.
\item
  If \(n = 0\) , A is a field. If \(n = 1\) , show that A is a principal ideal ring * and therefore a discrete valuation ring * (show that A is an integral domain by observing that \(m \cap m^n = \bigcap m^n\) , whence \(\bigcap m^n = 0\) ).
\item
  Show that the ring A{[}{[}T{]}{]} is regular.
\end{enumerate}

\begin{enumerate}
\def\labelenumi{\arabic{enumi})}
\setcounter{enumi}{13}
\item
  Let A be a ring, and let a be an ideal of A. Consider the graded alternating (A/a)-algebra \(\mathrm{Tor}^{\mathbf{A}}(\mathbf{A}/\mathbf{a}, \mathbf{A}/\mathbf{a})\) (cf. X, p. 201, exercise 9) and the homomorphism of graded (A/a)-algebras \(\theta: \Lambda_{\mathbf{A}/\mathbf{a}}(\mathbf{a}/\mathbf{a}^{2}) \to \mathrm{Tor}^{\mathbf{A}}(\mathbf{A}/\mathbf{a}, \mathbf{A}/\mathbf{a})\) deduced from the isomorphism \(a/a^{2} \to \mathrm{Tor}_{1}^{\mathbf{A}}(\mathbf{A}/\mathbf{a}, \mathbf{A}/\mathbf{a})\) (X, p. 73). Show that if the ideal a is generated by a completely secant sequence in A, \(\theta\) is an isomorphism.
\item
  Let A be a noetherian ring, and let a be an ideal contained in the radical of A. Suppose that the (A/a)-module \(a/a^{2}\) is free and that the homomorphism \(\theta_{2}: \Lambda_{A/a}^{2}(a/a^{2}) \to \operatorname{Tor}_{2}^{\mathbf{A}}(\mathbf{A}/\mathbf{a}, \mathbf{A}/\mathbf{a})\) (exercise 14) is surjective. Show that a is generated by a completely secant sequence in A. (If x is a minimal family of generators of a, show with the aid of exercise 7 (or directly) that Coker ( \(\theta_{2}\) ) is isomorphic to
\end{enumerate}

\[
\mathrm{H} _ {1} (\boldsymbol {x}, \mathbf {A}) \otimes_ {\mathbf {A}} \mathbf {A} / \mathfrak {a}.
\]

\begin{enumerate}
\def\labelenumi{\arabic{enumi})}
\setcounter{enumi}{15}
\tightlist
\item
  Let A be a local noetherian ring, m its maximal ideal, k = A/m. Show that the homomorphism of graded k-algebras \(\theta : \symbf{\Lambda}_{k}(\mathfrak{m}/\mathfrak{m}^{2}) \to \operatorname{Tor}^{\mathbf{A}}(k, k)\) is injective: if \(\theta_{2}\) is bijective, the ring A is regular (exercise 12).
\end{enumerate}

(If x is a minimal system of generators of m, show that the complex K(x, A) satisfies the conditions of exercise 14, p. 182; use exercises 4 and 15.)

¶ 17) Let A be a ring such that the A-module \(A_{s}\) has finite injective dimension (X, p. 202, exercise 4; a regular local ring (exercise 12) satisfies this property).

\begin{enumerate}
\def\labelenumi{\alph{enumi})}
\tightlist
\item
  Let x be a non-zero-divisor and non-invertible element of A. Show that one has
\end{enumerate}

\[
\mathrm{di} _ {\mathbf {A} / \mathbf {A} x} (\mathbf {A} / \mathbf {A} x) = \mathrm{di} _ {\mathbf {A}} (\mathbf {A}) - 1.
\]

(Let 0 → A → I⁰ → I¹ → \ldots{} be a minimal injective resolution (X, p. 182, exercise 13) of A; show that the complex Homₐ(A/AX, I¹) → Homₐ(A/AX, I²) → \ldots{} defines a minimal injective resolution of the (A/AX)-module A/AX.)

\begin{enumerate}
\def\labelenumi{\alph{enumi})}
\setcounter{enumi}{1}
\item
  We henceforth suppose that the ring A is local noetherian. Prove the equality di(A) = prof(A) (exercise 11; reason by induction on di(A)).
\item
  Let \((x_{1},\ldots ,x_{n})\) be a completely secant sequence for A, contained in the maximal ideal of A, with \(n = \mathrm{di}(\mathbf{A})\) . Show that the ring A is self-injective (X, p. 171, exercise 26).
\end{enumerate}

\begin{enumerate}
\def\labelenumi{\arabic{enumi})}
\setcounter{enumi}{17}
\tightlist
\item
  Let A be a noetherian ring, C a stable set of classes of A-modules (X, p. 40), \(x = (x_{1}, \ldots, x_{n})\)
\end{enumerate}

a sequence of elements of A, x the ideal generated by x. We denote by \(\mathcal{C}_{x}\) the set of elements of \(\mathcal{C}\) which are annihilated by x. Let M be a finitely generated A-module.

\begin{enumerate}
\def\labelenumi{\alph{enumi})}
\tightlist
\item
  Prove that for every \(i \geqslant 0\) , there exists an element \(\mathbf{M}_{i}\) of \(\mathcal{C}_{x}\) such that one has
\end{enumerate}

\[
\sum_ {k \geqslant 0} (- 1) ^ {k} \left[ \mathrm{H} _ {i + k} (\mathbf {x}, \mathrm{M}) \right] = [ \mathrm{M} _ {i} ]
\]

in \(\mathbf{K}(\mathcal{C}_{\pmb{x}})\) (reason by induction on \(n\) , using prop. 4, p. 157; in the case \(n = 1\) , consider the successive powers of the endomorphism \((x_{1})_{\mathbb{M}}\) ).

\begin{enumerate}
\def\labelenumi{\alph{enumi})}
\setcounter{enumi}{1}
\tightlist
\item
  We assume that \(\mathbf{A} / x\) has finite length. Show that there exist positive integers \(m_{i}\) ( \(i \geqslant 0\) ) such that one has length \(\mathrm{H}_{i}(x, \mathbf{M}) = m_{i} + m_{i+1}\) .
\end{enumerate}

\specialchapter{Index of Notations}

\[
(\mathbf {C}, d): \mathbf {p}. 2 3.
\]

\[
\mathrm{C} _ {n}, \mathrm{C} ^ {n}: \mathrm{p}. 2 4.
\]

\[
\mathbf {Z} (\mathbf {C}, d), \mathbf {Z} _ {n} (\mathbf {C}, d), \mathbf {B} (\mathbf {C}, d), \mathbf {B} _ {n} (\mathbf {C}, d), \mathbf {H} (\mathbf {C}),
\]

\[
\mathrm{H} _ {n} (\mathrm{C}, d), \mathrm{H} ^ {n} (\mathrm{C}, d): \mathrm{p}. 2 5.
\]

\[
\mathbf {H} (u), \mathbf {Z} (u), \mathbf {B} (u): \mathrm{p.26}.
\]

\[
\mathbf {C} (p): \mathrm{p.26}.
\]

\[
\mathbf {A} (\varepsilon): \mathbf {p}. 2 7.
\]

\[
\partial_ {u, v}, \partial : \mathbf {p}. 2 9.
\]

\[
\operatorname{Con} (u), \operatorname{Cyl} (u): \mathrm{p.36}.
\]

\[
\mathbf {K} (\mathcal {C}), [ \mathrm{M} ] _ {\mathcal {C}}, [ \mathrm{M} ]: \text { p. } 4 0.
\]

\[
\chi_ {\varphi} (\mathbf {M}), \chi (\mathbf {M}): \mathrm{p.~41}.
\]

\[
\Omega_ {\mathrm{A} / k} ^ {p}: \mathrm{p.~43.}
\]

\[
\nabla : \mathrm{p}. 4 5.
\]

\[
\mathbf {L} (\mathbf {M}), \mathbf {L} (f): \mathrm{p.} 5 0.
\]

\[
\mathrm{I} (\mathbf {M}), \mathrm{I} (f): \mathrm{p}. 5 2.
\]

\[
\mathbf {B} (\mathbf {A}, \mathbf {M}): \mathrm{p.} 5 7.
\]

\[
\mathbf {C} \otimes_ {\mathbf {A}} \mathbf {C} ^ {\prime}: \mathrm{p.~61}.
\]

\[
\begin{array}{l}\gamma (\mathbf {C}, \mathbf {C} ^ {\prime}): \mathbf {H} (\mathbf {C}) \otimes_ {\mathbf {A}} \mathbf {H} (\mathbf {C} ^ {\prime}) \rightarrow \mathbf {H} (\mathbf {C} \otimes_ {\mathbf {A}} \mathbf {C} ^ {\prime}):\\\text { p.62. }\end{array}
\]

\[
\sigma (\mathbf {C}, \mathbf {C} ^ {\prime}): \mathrm{p.~62.}
\]

\[
\operatorname{Tor} _ {n} ^ {\mathbf {A}} (\mathbf {M}, \mathbf {N}): \mathrm{p.~67}.
\]

\[
\operatorname{Tor} ^ {\mathbf {A}} (f, g): \mathrm{p.~68}.
\]

\[
\psi_ {\mathbf {M}} (\mathbf {N}), \overline {{\psi}} _ {\mathbf {N}} (\mathbf {M}): \mathrm{p.~69}.
\]

\[
\sigma_ {\mathbf {M}, \mathbf {N}}: \mathbf {p}. 7 1.
\]

\[
\operatorname{Homgr} _ {\mathbf {A}} (\mathbf {C}, \mathbf {C} ^ {\prime}): \text { p. } 8 1.
\]

\[
\lambda (\mathbf {C}, \mathbf {C} ^ {\prime}): \mathrm{p.} 8 2.
\]

\[
\operatorname{Ext} _ {\mathbf {A}} (\mathbf {M}, \mathbf {N}), \operatorname{Ext} _ {\mathbf {A}} ^ {n} (\mathbf {M}, \mathbf {N}): p. 8 6.
\]

\[
\varphi_ {\mathbf {M}} (\mathbf {N}), \overline {{\varphi}} _ {\mathbf {N}} (\mathbf {M}): \mathrm{p.} 8 8.
\]

\[
\delta (\mathbf {M}, \mathcal {E}): \mathrm{p.} 9 0.
\]

\[
\delta (\mathcal {F}, \mathrm{N}): \mathrm{p.} 9 2.
\]

\[
\psi (\mathrm{S}, \mathrm{R}), \varphi (\mathrm{R}, \mathrm{E}): \mathrm{p}. 1 0 0.
\]

\[
\psi^ {\prime} \left(\mathrm{S} _ {1}, \mathrm{R} _ {1}\right), \varphi^ {\prime} \left(\mathrm{R} _ {1}, \mathrm{E} _ {1}\right): \mathrm{p}. 1 0 2.
\]

\[
\mathrm{H} ^ {n} (\mathrm{G}, \mathrm{M}), \mathrm{H} _ {n} (\mathrm{G}, \mathrm{M}): \mathrm{p}. 1 1 1.
\]

\[
a _ {\mathbf {M}, \mathbf {N}}: \mathrm{p.} 1 1 3.
\]

\[
c _ {\mathbf {M}, \mathbf {N}, \mathbf {P}}, u \circ v: \mathrm{p}. 1 1 4.
\]

\[
\delta^ {m} (\mathbf {P}, \mathcal {S}), \delta^ {m} (\mathcal {S}, \mathbf {P}), \partial^ {m} (\mathbf {P}, \mathcal {S}), \partial^ {m} (\mathcal {S} _ {1}, \mathbf {M})
\]

\[
1 2 7, 1 3 1, 1 3 2.
\]

\[
c _ {\mathbf {P}, \mathbf {Q}; \mathbf {M}}: \mathbf {p}. 1 2 8.
\]

\[
c _ {\mathbf {P}; \mathbf {M}, \mathbf {N}}: \mathfrak {p}. 1 2 9.
\]

\[
\mathrm{dp} _ {\mathbf {A}} (\mathbf {M}): \mathrm{p}. 1 3 4.
\]

\[
\mathrm{dh} (\mathrm{A}): \mathrm{p}. 1 3 8.
\]

\[
\mathbf {K} (u), \mathbf {K} _ {\cdot} (u, \mathbf {C}), \mathbf {K} ^ {\prime} (u, \mathbf {C}): \mathrm{p.~147.}
\]

\[
b \vee b ^ {\prime}, a \top a ^ {\prime}: \mathrm{p.} 2 0 0, \text {   exerc.   } 7.
\]

\[
\mathrm{di} _ {\mathbf {A}} (\mathbf {M}): \mathbf {p}. 2 0 2, \text {   exerc.   } 4.
\]

\[
\mathrm{dpl} _ {\mathbf {A}} (\mathbf {M}): \text { p. } 2 0 2, \text { exerc. } 7.
\]

\[
\operatorname{td} (\mathbf {A}): \mathrm{p.202,exerc.8.}
\]

\[
\mathrm{da} _ {k} (\mathrm{A}): \mathrm{p.205,exerc.25.}
\]

\[
\operatorname{dc} (\mathbf {G}): \mathrm{p.206,exerc.28.}
\]

\specialchapter{Terminological Index}

Absolutely flat (ring ---): p. 169, exerc. 16. Acyclic in degree n (complex ---): p. 26. x-adic (topology ---): p. 159. Augmented algebra: p. 196, exerc. 27. Differential graded algebra: p. 183, exerc. 18. Ascending (degree ---): p. 23. Associated (class ---) to an exact sequence: p. 117. Augmentation of an algebra: p. 196, exerc. 27. Self-injective (ring): p. 171, exerc. 26. Bicomplex: p. 174, exerc. 11. Bounded above, below: p. 24. Canonical (injective resolution ---): p. 52. Canonical (free resolution ---): p. 50. Euler-Poincaré characteristic: p. 40. Cellular (complexes ---): p. 32. Golod-Shafarevich (theorem of ---): p. 192, exerc. 16. Five (lemma of the ---): p. 7. Universal coefficients (formula of ---): p. 78, 99. Cogenerator (module ---): p. 18. Coherent on the left, on the right (ring ---): p. 181, exerc. 11. Coherent, pseudo-coherent (module ---): p. 181, exerc. 10. Cohomology of groups: p. 111. Cohomological (dimension ---): p. 206, exerc. 28. Cohomologous: p. 194, exerc. 23. Co-induced (module ---): p. 190, exerc. 12. Commutation (isomorphism of ---): p. 71. Completely secant (family ---): p. 157. Complexes: p. 23. Koszul complexes: p. 147, 154. Composition (product of ---): p. 113. Cone of a morphism of complexes: p. 36. Conical (simplicial scheme ---): p. 177, exerc. 19. Connection: p. 45. Convergent (spectral sequence ---): p. 175, exerc. 14. Coregular (sequence ---): p. 159. Cylinder of a morphism of complexes: p. 36. Dedekind (rings of ---): p. 141, 203, exerc. 12. De Rham (complex of ---): p. 43. Descending (degree ---): p. 23. Commutative diagrams: p. 1. Differential: p. 24. Differential (algebra ---): p. 183, exerc. 18.

\begin{Shaded}
\begin{Highlighting}[]
\NormalTok{Exterior differential: p. 44.}
\NormalTok{Cohomological dimension: p. 206, exerc. 28.}
\NormalTok{Homological dimension: p. 138.}
\NormalTok{Injective dimension: p. 202, exerc. 4.}
\NormalTok{Flat dimension: p. 202, exerc. 7.}
\NormalTok{Projective dimension: p. 134.}
\NormalTok{Divisible (module ---): p. 17.}
\NormalTok{Dualizing (module ---): p. 172, exerc. 30.}

\NormalTok{Injective envelope: p. 19.}
\NormalTok{Euler{-}Poincaré (characteristic of\textquotesingle{}---): p. 40.}
\NormalTok{Extensions (modules of\textquotesingle{}---): p. 81, 86.}
\NormalTok{Extension of\textquotesingle{}an algebra: p. 197, exerc. 1.}
\NormalTok{Extension of\textquotesingle{}a Lie algebra: p. 198, exerc. 4.}

\NormalTok{Faithfully flat (module ---): p. 169, exerc. 9.}
\NormalTok{Differential forms: p. 44.}

\NormalTok{Goldie (theorem of —): p. 171, exerc. 25.}
\NormalTok{Golod{-}Chafarévitch (theorem of —): p. 192, exerc. 16.}

\NormalTok{Hilbert (syzygy theorem of —): p. 146.}
\NormalTok{Homology (classes of\textquotesingle{}—, module of\textquotesingle{}—) : p. 25.}
\NormalTok{Homology (exact sequence of\textquotesingle{}—): p. 30.}
\NormalTok{Group homology: p. 111.}
\NormalTok{Homological (dimension —): p. 138.}
\NormalTok{Homologism: p. 26.}
\NormalTok{Homologous (cycles —): p. 25.}
\NormalTok{Homomorphisms (complexes of\textquotesingle{}—) : p. 81.}
\NormalTok{Homotopic to zero (complex —): p. 34.}
\NormalTok{Homotopies (morphisms —): p. 32.}
\NormalTok{Homotopic (simplicially —): p. 177, exerc. 19.}
\NormalTok{Homotopy: p. 32.}
\NormalTok{Homotopism: p. 33.}

\NormalTok{Induced (module —): p. 190, exerc. 12.}
\NormalTok{Injective (modules —): p. 15.}
\NormalTok{Injective (dimension —): p. 202, exerc. 4.}
\NormalTok{Iterates (connecting homomorphisms —): pp. 127, 131, 132.}

\NormalTok{Koszul (complex of —): pp. 147, 154.}
\NormalTok{Künneth (formula of —): p. 79.}

\NormalTok{Connecting (homomorphism of —): p. 29.}
\NormalTok{Connecting (homomorphism of —) of extension modules: pp. 90, 92.}
\NormalTok{Connecting (homomorphism of —) of torsion products: p. 72.}
\NormalTok{Length of\textquotesingle{}a complex: p. 24.}
\NormalTok{Length (of\textquotesingle{}a presentation): p. 180, exerc. 6.}

\NormalTok{Minimal (injective resolution —): p. 182, exerc. 13.}
\NormalTok{Minimal (projective resolution —): p. 54.}
\NormalTok{Morphism of complexes: p. 25.}

\NormalTok{Zero on the right, zero on the left: p. 24.}
\end{Highlighting}
\end{Shaded}

\begin{Shaded}
\begin{Highlighting}[]
\NormalTok{Flat (absolutely —): p. 169, exerc. 16.}
\NormalTok{Flat (faithfully —): p. 169, exerc. 9.}
\NormalTok{Flat (dimension —): p. 202, exerc. 7.}
\NormalTok{Flat (modules —): p. 8.}
\NormalTok{Poincaré (characteristic of\textquotesingle{}Euler{-} ---): p. 40.}
\NormalTok{Poincaré polynomial: p. 41.}
\NormalTok{Precomplex: p. 174, exerc. 13.}
\NormalTok{Presentation, finite presentation: p. 10.}
\NormalTok{Presentation of length n: p. 180, exerc. 6.}
\NormalTok{Composition product: p. 113,}
\NormalTok{Torsion product: pp. 61, 67.}
\NormalTok{Projective (dimension —): p. 134.}
\NormalTok{Prüferian (ring —): p. 203, exerc. 11.}
\NormalTok{Pseudo{-}coherent (module —): p. 181, exerc. 10.}

\NormalTok{Regular (local ring —): p. 208, exerc. 12.}
\NormalTok{Regular (sequence —): p. 158.}
\NormalTok{Regular (spectral sequence —): p. 175, exerc. 14.}
\NormalTok{Relatively projective, injective: p. 170, exerc. 19.}
\NormalTok{Resolutions: p. 48.}
\NormalTok{Canonical injective resolution: p. 52.}
\NormalTok{Canonical free resolution: p. 50.}
\NormalTok{Graded resolution: p. 56.}
\NormalTok{Minimal projective resolution: p. 54.}
\NormalTok{Projective, injective resolution of\textquotesingle{}a complex: p. 183, exerc. 17.}

\NormalTok{Simplicial scheme: p. 177, exerc. 19.}
\NormalTok{Split (completely —): p. 35.}
\NormalTok{Splitting: p. 35.}
\NormalTok{Secant (completely —): p. 157.}
\NormalTok{Snake (— diagram): p. 3.}
\NormalTok{Simplex: p. 178, exerc. 19.}
\NormalTok{Simplicial (— scheme): p. 177, exerc. 19.}
\NormalTok{Simplicially homotopic: p. 177, exerc. 19.}
\NormalTok{Singular (— homology): p. 31.}
\NormalTok{Spectral (— sequence): p. 175, exerc. 14.}
\NormalTok{Standard (— resolution): p. 57.}
\NormalTok{Stationary (— spectral sequence): p. 175, exerc. 14.}
\NormalTok{Exact sequence of\textquotesingle{}homology: p. 30.}
\NormalTok{Spectral sequence: p. 175, exerc. 14.}
\NormalTok{Syzygies (theorem of —): p. 146.}

\NormalTok{Tensor (— product) of two complexes: p. 61.}
\NormalTok{Torsion (— product): p. 61, 67.}
\NormalTok{Tor{-}dimension: p. 202, exerc. 8.}
\NormalTok{Topology x{-}adic: p. 159.}
\NormalTok{(p{-}th) Translate: p. 27.}
\NormalTok{Transfer: p. 191, exerc. 14.}
\NormalTok{Trivial (cohomologically —): p. 192, exerc. 18.}

\NormalTok{Universal (— coefficients): p. 78, 99.}
\end{Highlighting}
\end{Shaded}
